\input eplain
\pdfpagewidth=148mm
\pdfpageheight=210mm
\hsize=128mm
\vsize=190mm

\pdfhorigin=0pt
\pdfvorigin=0pt
\hoffset=10mm
\voffset=10mm
\noindent gesture \leaders\hbox to 10pt {\hss.\hss}  \hfill \xref{node27}
\smallskip
\noindent bttyp \leaders\hbox to 10pt {\hss.\hss}  \hfill \xref{node36}
\smallskip
\noindent autoscaling \leaders\hbox to 10pt {\hss.\hss}  \hfill \xref{node19}
\smallskip
\noindent duplicate node \leaders\hbox to 10pt {\hss.\hss}  \hfill \xref{node22}
\smallskip
\noindent eps \leaders\hbox to 10pt {\hss.\hss}  \hfill \xref{node25}
\smallskip
\noindent fig prototype \leaders\hbox to 10pt {\hss.\hss}  \hfill \xref{node20}
\smallskip
\noindent initial parser \leaders\hbox to 10pt {\hss.\hss}  \hfill \xref{node18}
\smallskip
\noindent labels \leaders\hbox to 10pt {\hss.\hss}  \hfill \xref{node21}
\smallskip
\noindent notes integration \leaders\hbox to 10pt {\hss.\hss}  \hfill \xref{node24}
\smallskip
\noindent path drawing \leaders\hbox to 10pt {\hss.\hss}  \hfill \xref{node26}
\smallskip
\noindent sonilo \leaders\hbox to 10pt {\hss.\hss}  \hfill \xref{node0}
\smallskip
\noindent ink \leaders\hbox to 10pt {\hss.\hss}  \hfill \xref{node17}
\smallskip
\noindent news \leaders\hbox to 10pt {\hss.\hss}  \hfill \xref{node29}
\smallskip
\noindent tasks \leaders\hbox to 10pt {\hss.\hss}  \hfill \xref{node1}
\smallskip
\noindent allocator \leaders\hbox to 10pt {\hss.\hss}  \hfill \xref{node12}
\smallskip
\noindent block pool \leaders\hbox to 10pt {\hss.\hss}  \hfill \xref{node37}
\smallskip
\noindent blockstack improvements \leaders\hbox to 10pt {\hss.\hss}  \hfill \xref{node16}
\smallskip
\noindent blsaw \leaders\hbox to 10pt {\hss.\hss}  \hfill \xref{node10}
\smallskip
\noindent buddy slots \leaders\hbox to 10pt {\hss.\hss}  \hfill \xref{node14}
\smallskip
\noindent buddy system \leaders\hbox to 10pt {\hss.\hss}  \hfill \xref{node3}
\smallskip
\noindent bulk instructions \leaders\hbox to 10pt {\hss.\hss}  \hfill \xref{node31}
\smallskip
\noindent butlp \leaders\hbox to 10pt {\hss.\hss}  \hfill \xref{node23}
\smallskip
\noindent docgen \leaders\hbox to 10pt {\hss.\hss}  \hfill \xref{node13}
\smallskip
\noindent fix block navigation \leaders\hbox to 10pt {\hss.\hss}  \hfill \xref{node28}
\smallskip
\noindent ilo interpreter \leaders\hbox to 10pt {\hss.\hss}  \hfill \xref{node4}
\smallskip
\noindent initial dsp \leaders\hbox to 10pt {\hss.\hss}  \hfill \xref{node9}
\smallskip
\noindent instruction primitive \leaders\hbox to 10pt {\hss.\hss}  \hfill \xref{node30}
\smallskip
\noindent konilo sonilo communication \leaders\hbox to 10pt {\hss.\hss}  \hfill \xref{node5}
\smallskip
\noindent meaningful sounds \leaders\hbox to 10pt {\hss.\hss}  \hfill \xref{node2}
\smallskip
\noindent mem \leaders\hbox to 10pt {\hss.\hss}  \hfill \xref{node6}
\smallskip
\noindent phasor \leaders\hbox to 10pt {\hss.\hss}  \hfill \xref{node8}
\smallskip
\noindent sonilo command language \leaders\hbox to 10pt {\hss.\hss}  \hfill \xref{node7}
\smallskip
\noindent ugen bytecode \leaders\hbox to 10pt {\hss.\hss}  \hfill \xref{node11}
\smallskip
\noindent ugen ports \leaders\hbox to 10pt {\hss.\hss}  \hfill \xref{node33}
\smallskip
\noindent ugen refcount \leaders\hbox to 10pt {\hss.\hss}  \hfill \xref{node34}
\smallskip
\noindent ugen stack \leaders\hbox to 10pt {\hss.\hss}  \hfill \xref{node35}
\smallskip
\noindent ugens \leaders\hbox to 10pt {\hss.\hss}  \hfill \xref{node32}
\smallskip
\noindent zero page \leaders\hbox to 10pt {\hss.\hss}  \hfill \xref{node15}
\smallskip
\vfill \break
{\medskip\noindent\bf\font\titlefont=cmss17 at 20pt\titlefont Timeline}\medskip
{\noindent \bf 2025-10-22:} sonilo, tasks, meaningful sounds\smallskip
{\noindent \bf 2025-10-29:} buddy system\smallskip
{\noindent \bf 2025-11-01:} buddy system\smallskip
{\noindent \bf 2025-11-02:} buddy system, ilo interpreter\smallskip
{\noindent \bf 2025-11-05:} buddy system\smallskip
{\noindent \bf 2025-11-08:} buddy system\smallskip
{\noindent \bf 2025-11-09:} buddy system, konilo sonilo communication\smallskip
{\noindent \bf 2025-11-11:} konilo sonilo communication\smallskip
{\noindent \bf 2025-11-13:} buddy system, ilo interpreter, konilo sonilo communication, mem, phasor, sonilo command language\smallskip
{\noindent \bf 2025-11-14:} ilo interpreter\smallskip
{\noindent \bf 2025-11-15:} buddy system\smallskip
{\noindent \bf 2025-11-16:} buddy system\smallskip
{\noindent \bf 2025-11-24:} ilo interpreter, initial dsp, phasor\smallskip
{\noindent \bf 2025-11-25:} phasor\smallskip
{\noindent \bf 2025-11-26:} ilo interpreter\smallskip
{\noindent \bf 2025-11-28:} sonilo, ilo interpreter, mem\smallskip
{\noindent \bf 2025-11-29:} blsaw\smallskip
{\noindent \bf 2025-11-30:} tasks, allocator, docgen, ilo interpreter, ugen bytecode\smallskip
{\noindent \bf 2025-12-03:} sonilo, buddy slots\smallskip
{\noindent \bf 2025-12-12:} buddy slots\smallskip
{\noindent \bf 2025-12-13:} buddy slots\smallskip
{\noindent \bf 2025-12-16:} buddy slots, zero page\smallskip
{\noindent \bf 2025-12-17:} zero page\smallskip
{\noindent \bf 2025-12-18:} blockstack improvements, zero page\smallskip
{\noindent \bf 2025-12-19:} blockstack improvements\smallskip
{\noindent \bf 2025-12-21:} ink\smallskip
{\noindent \bf 2025-12-28:} initial parser\smallskip
{\noindent \bf 2025-12-29:} autoscaling, fig prototype, initial parser\smallskip
{\noindent \bf 2025-12-30:} labels\smallskip
{\noindent \bf 2025-12-31:} duplicate node\smallskip
{\noindent \bf 2026-01-02:} notes integration, tasks, allocator, blockstack improvements, buddy slots, butlp, ilo interpreter, konilo sonilo communication, sonilo command language, zero page\smallskip
{\noindent \bf 2026-01-04:} blockstack improvements\smallskip
{\noindent \bf 2026-01-05:} blockstack improvements\smallskip
{\noindent \bf 2026-01-06:} notes integration, blockstack improvements\smallskip
{\noindent \bf 2026-01-07:} notes integration\smallskip
{\noindent \bf 2026-01-08:} blockstack improvements\smallskip
{\noindent \bf 2026-01-10:} eps, path drawing\smallskip
{\noindent \bf 2026-01-11:} gesture, notes integration, blockstack improvements\smallskip
{\noindent \bf 2026-01-14:} sonilo, blockstack improvements\smallskip
{\noindent \bf 2026-01-15:} blockstack improvements, buddy slots, fix block navigation, zero page\smallskip
{\noindent \bf 2026-01-16:} blockstack improvements, buddy slots, fix block navigation, sonilo command language\smallskip
{\noindent \bf 2026-01-17:} buddy slots\smallskip
{\noindent \bf 2026-01-18:} ugen bytecode\smallskip
{\noindent \bf 2026-01-22:} news\smallskip
{\noindent \bf 2026-01-24:} sonilo, bulk instructions, instruction primitive, ugen bytecode, ugens\smallskip
{\noindent \bf 2026-01-26:} ugens\smallskip
{\noindent \bf 2026-01-31:} bttyp, bulk instructions, ilo interpreter, sonilo command language, ugen ports, ugen refcount, ugen stack, ugens\smallskip
{\noindent \bf 2026-02-03:} bttyp, news\smallskip
{\noindent \bf 2026-02-07:} news, fix block navigation, ugen refcount\smallskip
{\noindent \bf 2026-02-14:} news, block pool, bulk instructions, ugen stack\smallskip
{\noindent \bf 2026-02-15:} news\smallskip
\vfill \break
\xrdef{node0}
{\medskip\noindent\bf\font\titlefont=cmss17 at 20pt\titlefont sonilo}\medskip
{\medskip\noindent\bf\font\titlefont=cmitt12 at 17pt\titlefont 2025-10-22}\medskip
{\smallskip\noindent\bf\font\titlefont=cmr12 at 13pt\titlefont $\bullet$ 14:04: Transcribing Task Graphs for Sonilo}\smallskip
It needs to be fleshed out more, but I want real data to work with for this task graphing utility I am going to build for scheduling tasks and generating daily agendas.
\par

I ended up building the high-level "artist purpose" knowledge graph I started building yesterday, which outline artistic purpose. The actual tasks have yet to be made.
\par

{\medskip\noindent\bf\font\titlefont=cmitt12 at 17pt\titlefont 2025-11-28}\medskip
{\smallskip\noindent\bf\font\titlefont=cmr12 at 13pt\titlefont $\bullet$ 09:17: addressing the agenda overflow}\smallskip
As of today, there are 26 tasks total.
\par

9 of those tasks are related to sonilo. I am stuck. Memory allocators have taken time. Ilo has taken time. Trying to get the moving parts right is taking time. Sonilo ilo interpreter is about 26 days overdue, that's the most overdue task.
\par

COAD also has 9 tasks, surprisingly. I guess this makes a little more sense. I'll have gone about 10 days without doing any reading. The most overdue task is 13 days, so I'm about 3-5 days overdue factoring in this trip.
\par

DMOI is 11 days overdue for the latest. That means I'm about 1-3 days overdue factoring vacation. Not too bad.
\par

I am going to add in a way to get the max day for selected groups, then shift sonilo and COAD back 10 days. Once I have this "max days" functionality, I should be able to calculate a "due date" from the agenda.
\par

{\medskip\noindent\bf\font\titlefont=cmitt12 at 17pt\titlefont 2025-12-03}\medskip
{\smallskip\noindent\bf\font\titlefont=cmr12 at 13pt\titlefont $\bullet$ 11:02: shifting goalposts again for sonilo project, dmoi}\smallskip
I need these tasks to go down, so I'm redistributing the days. I'm disappointed at velocity, but at least I'm measuring it. the number is now 180 days
\par

stretching the reading for dmoi so only 2 tasks appear right now the reading takes a while. I am a thorough (read: slow) reader.
\par

{\medskip\noindent\bf\font\titlefont=cmitt12 at 17pt\titlefont 2026-01-14}\medskip
{\smallskip\noindent\bf\font\titlefont=cmr12 at 13pt\titlefont $\bullet$ 10:23: I'm scared to look at the sonilo tasks}\smallskip
{\medskip\noindent\bf\font\titlefont=cmitt12 at 17pt\titlefont 2026-01-24}\medskip
{\smallskip\noindent\bf\font\titlefont=cmr12 at 13pt\titlefont $\bullet$ 09:43: sonilo saturdays}\smallskip
I'm trying to establish a "sonilo saturdays" habit. Where I dedicate a good chunk of every saturday developing sonilo
\par

{\smallskip\noindent\bf\font\titlefont=cmr12 at 13pt\titlefont $\bullet$ 09:44: logs review. where am I today?}\smallskip
Looks like I finished up the zero page stuff.
\par

I am working on methods to convert outputs from the block stack allocator (block id numbers) to word addresses in memory (b to m), and back again (m to b).
\par

I need a chronological view of things in my log output. Date, then what things were worked on on that day.
\par

There's an issue with my block to memory convertors in that the memory address returned isn't the smallest value. I still need to make an issue for it.
\par

In sonilo land, I was in the middle of building up a test for the 'sm' command or "slot make". This tool appends a slot to the buddy slot list.
\par

The big idea I had last week with regards to the buddy slot stuff was to create operations that could interop with the existing single-instance buddy system operations.
\par

I think I also had an idea about returning to the DSP bytecode. I think I have enough allocation stuff working that I can work on this.
\par

DSP ugen stuff might be a focus of mine because I'm losing momentum with memory allocators.
\par

\xrdef{node1}
{\medskip\noindent\bf\font\titlefont=cmss17 at 20pt\titlefont tasks}\medskip
{\medskip\noindent\bf\font\titlefont=cmitt12 at 17pt\titlefont 2025-10-22}\medskip
{\smallskip\noindent\bf\font\titlefont=cmr12 at 13pt\titlefont $\bullet$ 15:47: back to tasks for sonilo (hopefully)}\smallskip
{\smallskip\noindent\bf\font\titlefont=cmr12 at 13pt\titlefont $\bullet$ 16:40: finish the initial graph traversal program}\smallskip
Also used this time to clean up the task graph. I now have top level task in Sonilo that just reads "make meaningful sounds with sonilo".
\par

{\medskip\noindent\bf\font\titlefont=cmitt12 at 17pt\titlefont 2025-11-30}\medskip
{\smallskip\noindent\bf\font\titlefont=cmr12 at 13pt\titlefont $\bullet$ 09:02: reworking some of the tasks}\smallskip
{\medskip\noindent\bf\font\titlefont=cmitt12 at 17pt\titlefont 2026-01-02}\medskip
{\smallskip\noindent\bf\font\titlefont=cmr12 at 13pt\titlefont $\bullet$ 14:07: reflecting on tasks}\smallskip
{\smallskip\noindent\bf\font\titlefont=cmr12 at 13pt\titlefont $\bullet$ 15:21: more sonilo check-ins}\smallskip
\xrdef{node2}
{\medskip\noindent\bf\font\titlefont=cmss17 at 20pt\titlefont meaningful sounds}\medskip
{\medskip\noindent\bf\font\titlefont=cmitt12 at 17pt\titlefont 2025-10-22}\medskip
{\smallskip\noindent\bf\font\titlefont=cmr12 at 13pt\titlefont $\bullet$ 16:40: finish the initial graph traversal program}\smallskip
Also used this time to clean up the task graph. I now have top level task in Sonilo that just reads "make meaningful sounds with sonilo".
\par

\xrdef{node3}
{\medskip\noindent\bf\font\titlefont=cmss17 at 20pt\titlefont buddy system}\medskip
{\medskip\noindent\bf\font\titlefont=cmitt12 at 17pt\titlefont 2025-10-29}\medskip
{\smallskip\noindent\bf\font\titlefont=cmr12 at 13pt\titlefont $\bullet$ 11:28: Work on the buddy system some more}\smallskip
Also, trying on this task management system.
\par

{\smallskip\noindent\bf\font\titlefont=cmr12 at 13pt\titlefont $\bullet$ 15:19: more buddy system}\smallskip
{\smallskip\noindent\bf\font\titlefont=cmr12 at 13pt\titlefont $\bullet$ 19:25: pre-buddy system, buddy system proto-type?}\smallskip
I am running into constraints due to the fact that the buddy system needs to be pre-initialized because it's modified to fit entirely on a single block. It would be easier to do this closer to the "regular" way, and then refactor into the block system.
\par

It was a good idea to build this, but it isn't working fully yet. Will need to troubleshoot at another time.
\par

{\medskip\noindent\bf\font\titlefont=cmitt12 at 17pt\titlefont 2025-11-01}\medskip
{\smallskip\noindent\bf\font\titlefont=cmr12 at 13pt\titlefont $\bullet$ 07:29: pre-coffee puttering, hack on buddy system}\smallskip
{\smallskip\noindent\bf\font\titlefont=cmr12 at 13pt\titlefont $\bullet$ 11:29: more buddy system}\smallskip
I think alloc() works now? Two mistakes made: had to be more explicit with what LOC(AVAIL[j]) was doing. I was using virtual pointers for everything, which made things interesting. There was also (x<-y<-z) notation used, which I thought would be intermediate assignments like "x = y, y = z", but I think it meant "set x = y = z".
\par

{\smallskip\noindent\bf\font\titlefont=cmr12 at 13pt\titlefont $\bullet$ 15:46: onto block "liberation"}\smallskip
For the record, the last bug was found by claude. It was so subtle, only 1 line, I transcribed it wrong, and knowing me it probably saved me about an hour of debugging and careful checking.
\par

{\smallskip\noindent\bf\font\titlefont=cmr12 at 13pt\titlefont $\bullet$ 17:23: what's next for this}\smallskip
I gotta adapt the algorithm to work with a 64-sized block. And, I gotta remove the gotos.
\par

I've learned a few things writing this out. I've got things set up in a way the algorithm is broken down into core operations that I can re-implement without too much trouble. Then, I can re-use the rest of the bits.
\par

I need a way to bootstrap the initial block, and permanently allocate a region of size 1 for the avail list, and a region of size 2 for the tags. Those should predictably appear in the first 3 words of the block. The temporary memory storying locations for the avail list and the tags can then be copied over to those memory locations, and the block can be used as-is for the remainder of the lifetime.
\par

The avail list has quirky behavior. when m = 6 it's a bool, otherwise, it's an actual location. This logic needs to be abstracted away cleanly.
\par

{\smallskip\noindent\bf\font\titlefont=cmr12 at 13pt\titlefont $\bullet$ 20:57: buddy system initial refactoring}\smallskip
I need to get this ready for my block system approach.
\par

Some initial refactoring is done. I'm trying to pull out the tag and AVAIL list logic. I'm now at the point where I'm getting the AVAIL list reworked so it can be external. I think I've run into a small problem: I've assumed that the AVAIL list can fit into 32 bits, and I don't think that's true. I think each entry needs to be at least 12 bits, not 6 bits because it needs LINKF/LINKB addresses. Also, I think I need 7 slots, not 6. So that's 84 bits, which can fit onto 3 words. So that's 5 words total instead of 3. So that's 58 words available, ~90% capacity self-contained.
\par

I am hoping we don't need kval for this, otherwise that's another word required. Each entry becomes 15 bits of information, and that's 105 bits.
\par

Alternatively, the state could be managed on another temp block, which can be returned after the memory is set. Ideally, I'll have enough flexibility to support both.
\par

{\medskip\noindent\bf\font\titlefont=cmitt12 at 17pt\titlefont 2025-11-02}\medskip
{\smallskip\noindent\bf\font\titlefont=cmr12 at 13pt\titlefont $\bullet$ 18:39: more of this buddy system}\smallskip
I flipped around the ordering of LINKF,LINKB,and KVAL so that LINKF and LINKB are in the lower bits. My hope is that the AVAIL table only needs those two values, and I can just re-use the linkf/linkb routines.
\par

{\smallskip\noindent\bf\font\titlefont=cmr12 at 13pt\titlefont $\bullet$ 20:15: back to it}\smallskip
It's a lot of work to cram this all in one block, especially considering that there will be some kind of wasted space. Having a temporary block for bookkeeping seems like the best option.
\par

6 words will be needed (rounded up to 8, who knows this temp block could be broken off from another buddy system).
\par

4 words are needed for the AVAIL list. the LINKF/LINKB entries will be pointers for a 6-bit address space with an extra bit to indicate addresses on the temporary block. Each entry is 14 bits, that's 2 bits a word with padding, 4 words to get 7 entries.
\par

tags will need 64 bits, so that's the remaining 2 words.
\par

I've created availf/b v2 functions. They seem to work in isolation, but when I put them in the program, it causes an infinite loop to happen when the checksums. I'm guessing it's some kind of addressing issue. I'll have to give it a look another time.
\par

{\medskip\noindent\bf\font\titlefont=cmitt12 at 17pt\titlefont 2025-11-05}\medskip
{\smallskip\noindent\bf\font\titlefont=cmr12 at 13pt\titlefont $\bullet$ 09:50: more buddy system}\smallskip
Last time, I got some initial availf/availb v2 logic coded up, which allows the avail list to be external. However, when I coded it up, the program got caught in a loop. I think it was where the traversal logic happened.
\par

The checksums are wrong as well, so I'm overlooking something.
\par

I'm writing a test program for AVAILF/AVAILB to see if there are differents in output.
\par

It looks like AVAILF is not being set properly. It doesn't seem to be being set at all. It is being initialized correctly. It's only after allocating a block that the AVAILF fields start to misalign.
\par

To be continued...
\par

{\smallskip\noindent\bf\font\titlefont=cmr12 at 13pt\titlefont $\bullet$ 16:55: AVAIL investigations continued}\smallskip
In the first allocation AVAILF(k) is supposed to be set to be 70, but it never seems to be set, staying at 0. The availf never explicitly sets itself to be 70. So linkf is making that call I think.
\par

I'm starting to see the issues. This new system needs to create an abstraction for converting virtual memory addresses into actual memory addresses. This algorithm currently assumes AVAIL is contiguously after the block (0-63 are the block address, 64-70 are the AVAIL entries).
\par

I'm thinking of going back to one word to one entry. Space is not an issue anymore because it's only temporarily allocated. Then, there needs to be some kind of virtual to physical memory resolver. So, anything that is between 0-63 gets an address from the block, anything greater than that gets an address from the temp block. I think if we do that, it's going to get much simpler.
\par

{\medskip\noindent\bf\font\titlefont=cmitt12 at 17pt\titlefont 2025-11-08}\medskip
{\smallskip\noindent\bf\font\titlefont=cmr12 at 13pt\titlefont $\bullet$ 21:43: yup, buddy system}\smallskip
I got the v2 system to match! Now to see if it lines up with the first test I made.
\par

Yup, it works now.
\par

What is going to happen next is, I need to refactor this to use an external block of memory. If TAGS and AVAIL are lumped together, that's 9 words in a 16 word block (rounded up). If TAGS and AVAIL are two separate blocks, that's 2 words (TAGS, 2 words), and 8 words (AVAIL, 7 words). So, I'm going to allocate AVAIL and TAGS separately in my buddy system.
\par

I should probably also simulate the memory of sonilo by using 2 blocks (2*64 = 128 words) of pre-allocated memory.
\par

progress is slow, but it is still progress.
\par

{\medskip\noindent\bf\font\titlefont=cmitt12 at 17pt\titlefont 2025-11-09}\medskip
{\smallskip\noindent\bf\font\titlefont=cmr12 at 13pt\titlefont $\bullet$ 18:02: refactoring sonilo buddy system to use virtual pointers}\smallskip
This will make it one step closer towards being compatible with the Sonilo memory model.
\par

So close to being done. I need the pointers themselves to be in the memory. From there, I'll need to be able to create new functions that work with the chunk of memory directly.
\par

{\smallskip\noindent\bf\font\titlefont=cmr12 at 13pt\titlefont $\bullet$ 18:26: wrap in the virtual pointers into the memory block}\smallskip
To be continued...
\par

{\smallskip\noindent\bf\font\titlefont=cmr12 at 13pt\titlefont $\bullet$ 21:00: ...continued}\smallskip
{\medskip\noindent\bf\font\titlefont=cmitt12 at 17pt\titlefont 2025-11-13}\medskip
{\smallskip\noindent\bf\font\titlefont=cmr12 at 13pt\titlefont $\bullet$ 10:20: buddy system: almost done}\smallskip
I think this is almost done. The algorithm works, the initial API has been written for the most part. I'd like to get this working in my growing test suite.
\par

{\medskip\noindent\bf\font\titlefont=cmitt12 at 17pt\titlefont 2025-11-15}\medskip
{\smallskip\noindent\bf\font\titlefont=cmr12 at 13pt\titlefont $\bullet$ 09:09: work towards wrapping up this buddy system lol}\smallskip
oh, right. this blocker is a bit of a ridiculous situation. I'm trying to get this into my test program, by my test program isn't variadic. I need multiple arguments to initialize the buddy system, but I only have one I/O register, and an address register.
\par

time to yak shave a stack?
\par

Set up some test code for an array with append and pop operations. Will need to write out code later.
\par

I need to tell myself this will help me later. It is interesting getting into the groove of this little system I build for myself.
\par

{\smallskip\noindent\bf\font\titlefont=cmr12 at 13pt\titlefont $\bullet$ 13:35: implement the stack}\smallskip
{\smallskip\noindent\bf\font\titlefont=cmr12 at 13pt\titlefont $\bullet$ 14:40: now, use this stack in the tests for the memory system}\smallskip
{\smallskip\noindent\bf\font\titlefont=cmr12 at 13pt\titlefont $\bullet$ 15:56: continue setting up the tests}\smallskip
It's been slow going because I'm dogfooding this test system, and learning how to use it, filling in missing parts along the way.
\par

I'm hoping things will get a bit faster? I'm reaching a pretty good stride with ed. the interactive test system also works pretty well with gdb. I can run components of the program in a very adhoc way, and then use gdb to add breakpoints and step through the program.
\par

oh, also, had to do some non-trivial debugging of the system for the first time. The test suite with the diffed output and gdb were a great combo.
\par

{\smallskip\noindent\bf\font\titlefont=cmr12 at 13pt\titlefont $\bullet$ 21:52: more test scaffolding}\smallskip
The expected test output has been written out, now I need to implement the programs, and then the commands.
\par

{\medskip\noindent\bf\font\titlefont=cmitt12 at 17pt\titlefont 2025-11-16}\medskip
{\smallskip\noindent\bf\font\titlefont=cmr12 at 13pt\titlefont $\bullet$ 10:46: buddy implementation continued}\smallskip
Damn, I'm not storing these addresses anywhere, so freeing may be tricky for this test. I'd have to invent variables. No time for this today.
\par

Maybe I can the input address be the relative block offset? This guarantees that the address belongs to the block. Getting the offset from a word address is just a matter of calculating the offset from the word address of the start of the block.
\par

I could also just compute the address. I know the base address of memory, so I just add the offset to it. 'ad' could be used for 'addition'. and 'sb', 'di', and 'mu' could be reserved for subtract, multiply, divide.
\par

the checksummer is hanging again. I think my problem has to do with address space. I think the buddy allocator uses local offset addresses, not word memory addresses. I'm passing in full word addresses.
\par

So actually, I think the local offsets are the expected arguments?
\par

Works now, nice! I think we're done.
\par

{\smallskip\noindent\bf\font\titlefont=cmr12 at 13pt\titlefont $\bullet$ 12:17: this is completed. finally}\smallskip
\xrdef{node4}
{\medskip\noindent\bf\font\titlefont=cmss17 at 20pt\titlefont ilo interpreter}\medskip
{\medskip\noindent\bf\font\titlefont=cmitt12 at 17pt\titlefont 2025-11-02}\medskip
{\smallskip\noindent\bf\font\titlefont=cmr12 at 13pt\titlefont $\bullet$ 13:43: initial exploration}\smallskip
I think this is mainly about adapting the existing ilo byte code interpreter and pali assembler as much as possible.
\par

The VM code seems like it'll mostly do the trick. I'll probably want to simplify it a bit. I do not need or want disk I/O so the blocks file should be optional.
\par

pali seems to be a system mostly designed to produce a singular ilo.rom file. I'm wondering if I could adapt the TAL assembler to work with ILO opcodes instead of Uxn.
\par

hang on, my 1-line test program works with pali. I might be able to work with this in the very short-term.
\par

At some point, I'll need more of an API for reading ILO streams.
\par

Wrote a dumb little program to print a character and multiply some numbers, with some jump logic. Good enough for now.
\par

{\medskip\noindent\bf\font\titlefont=cmitt12 at 17pt\titlefont 2025-11-13}\medskip
{\smallskip\noindent\bf\font\titlefont=cmr12 at 13pt\titlefont $\bullet$ 10:21: ilo interpreter: try to port uxntal?}\smallskip
Today, I am hoping to put some time into studying the source code for uxntal. It's not many lines. I think it could be reworked to generate ilo roms instead of uxn roms.
\par

For the initial prototyping stage, I imagine I'll have a bunch of ad-hoc programs generating ilo bytecode. pali, the assembler used to generate ilo.rom, is really only designed to genenerate that one rom file. I need smaller programs.
\par

enscript is a neat little tool I learned about that turns textfiles into postscript. Postscript to PDF means I can put it on my RM. I'm hoping this workflow will help me better study code more deeply.
\par

{\smallskip\noindent\bf\font\titlefont=cmr12 at 13pt\titlefont $\bullet$ 10:45: printing uxntal to PDF}\smallskip
command: === enscript -o - uxnasm.c | ps2pdf - uxnasm.pdf ===
\par

{\medskip\noindent\bf\font\titlefont=cmitt12 at 17pt\titlefont 2025-11-14}\medskip
{\smallskip\noindent\bf\font\titlefont=cmr12 at 13pt\titlefont $\bullet$ 17:35: uxnasm studying}\smallskip
I'm thinking trying to work uxnasm to output ilo code of uxn code. To prep for this, I thought it would be a good exercise to try and build a graph around the program. I put a monochrome version of uxnasm.c on my remarkable, and started drawing a graph of the program using this new graph system I've been developing. It was very time consuming.
\par

{\smallskip\noindent\bf\font\titlefont=cmr12 at 13pt\titlefont $\bullet$ 21:25: wrapping up uxnasm knowledge graph}\smallskip
While I don't understand the details of this program, I have a pretty good sense of what all the components of the program are and where they are in the file. The "parse" function is the most heavily loaded. I also need to somehow rework "ops" as well.
\par

{\medskip\noindent\bf\font\titlefont=cmitt12 at 17pt\titlefont 2025-11-24}\medskip
{\smallskip\noindent\bf\font\titlefont=cmr12 at 13pt\titlefont $\bullet$ 11:27: writing out pseudo-code for an ILO program computing phasor}\smallskip
This exercise is exposing some points of friction for me.
\par

The hard part is memory structure, once again. The ILO bits are very straight forward, once the memory structure is set up.
\par

Point of friction: I started this task wanting to just write out pseudo-assembler code for ILO. Then I realized, I need memory addresses. How do I generate those things? I'd need to allocate and initialize data for the phasor struct, in addition to the parameter.
\par

So, to make a program is really a collaboration between the memory writer and the ilo assembler. I'd use the memory writer to set up the memory structure, and the assembler to generate ilo bytecode.
\par

Ideally, it would be nice for programs to use the original ilo C implementation to start. I can make it so the ilo program is at the start of the file, with all the special memory after it. The program would terminate by a goto that goes beyond the ilo bounds.
\par

ilo programs are going to be very small, and the address space localized to just a handful of blocks. This phasor program would exist on one block.
\par

What does this phasor program do? It opens a file for writing, runs a loop writing stuff to that file, then closes that file. In the loop, the phasor ugen is called to generate a block of signal. That signal then gets written to disk. It is assumed that the phasor data is written to disk pre-initialized.
\par

Okay, this is clearly more work than I expected.
\par

{\medskip\noindent\bf\font\titlefont=cmitt12 at 17pt\titlefont 2025-11-26}\medskip
{\smallskip\noindent\bf\font\titlefont=cmr12 at 13pt\titlefont $\bullet$ 09:45: today's ilo interpreter thoughts}\smallskip
Current difficulty: memory structure vs ilo VM. These are two separate systems, each with their own command space. How does one bridge the two systems?
\par

When I was writing pseudo-code for the phasor, there were two explicit steps: setting up memory, and then setting up the ILO instructions. The memory was set up using my ad-hoc memory language, which then resulted in a single top-level memory address. The ILO code would push a constant parameter to a stack, push the pointer to the DSP to the stack, then call the DSP.
\par

It's a little bumpy at the moment.
\par

Well, the family is calling. that's the show. Might bring pen and paper with me today to the house.
\par

{\medskip\noindent\bf\font\titlefont=cmitt12 at 17pt\titlefont 2025-11-28}\medskip
{\smallskip\noindent\bf\font\titlefont=cmr12 at 13pt\titlefont $\bullet$ 09:49: custom bytecode language for DSP?}\smallskip
I'm starting to think IO calls in ILO are going to have too much overhead for the ugen related things. I think it would be better to have a bytecode interpreter more specialized for the DSP code, optimized to fit on 64-word blocks.
\par

I'm imagining something similar to the command-packing format of the ILO VM, except that there'd only be 16 opcodes (4-bit, 8 commands/word). Most of the functionality needed is pushing constant values calling hard-coded subroutines.
\par

Memory initialization could also probably happen using this system as well. There will need to be memory primitives for allocating small chunks of memory, as well as allocating full blocks (to write audio to). DSP unit generators would then have an initialization function that could set up the data in memory.
\par

{\medskip\noindent\bf\font\titlefont=cmitt12 at 17pt\titlefont 2025-11-30}\medskip
{\smallskip\noindent\bf\font\titlefont=cmr12 at 13pt\titlefont $\bullet$ 08:59: ilo interpreter can wait, use custom bytecode format instead}\smallskip
I think the DSP is specialized enough that I can get away with it.
\par

Ilo, on the other hand, is going to be more important in situtaions where small, inline executable programs are desirable. I'm mainly thinking this could be used for  gesture iterators.
\par

{\medskip\noindent\bf\font\titlefont=cmitt12 at 17pt\titlefont 2026-01-02}\medskip
{\smallskip\noindent\bf\font\titlefont=cmr12 at 13pt\titlefont $\bullet$ 14:39: check-in: how useful is ilo?}\smallskip
Since I originally wrote this task, the role I had in mind for Ilo has changed. Initially, I was going to use it for constructing DSP chains, but I've been reconsidering that.
\par

I think ilo will become more useful for more highlevel structures like gesture. Though, to be honest, I haven't thought that far ahead beyond how it'd be used.
\par

{\medskip\noindent\bf\font\titlefont=cmitt12 at 17pt\titlefont 2026-01-31}\medskip
{\smallskip\noindent\bf\font\titlefont=cmr12 at 13pt\titlefont $\bullet$ 16:01: task ordering updates, ugens child of some other tasks now}\smallskip
ilo interpreter stuff needs to at least happen after ugen stuff. And it's a prereq for any gesture stuff.
\par

By the time unit generators are finished, I think thinking about using a higher-level abstraction like Konilo may be warranted. At that point, it'll start to feel possible to build demos.
\par

\xrdef{node5}
{\medskip\noindent\bf\font\titlefont=cmss17 at 20pt\titlefont konilo sonilo communication}\medskip
{\medskip\noindent\bf\font\titlefont=cmitt12 at 17pt\titlefont 2025-11-09}\medskip
{\smallskip\noindent\bf\font\titlefont=cmr12 at 13pt\titlefont $\bullet$ 10:21: quick words about konilo-sonilo communication}\smallskip
This has popped up in my agenda, but I am behind because of the buddy memory allocator. To keep momentum going, I'm going to write down what I'm thinking about today.
\par

I'm not 100% sold on using Konilo. I am a bit more committed to using Ilo.
\par

Communication, sonilo or otherwise, is going to boil down to being a serial communication format. In Konilo, this would be done using IO commands.
\par

My initial thinking with using Konilo was that since both systems used Ilo, there could be some convenience there. It also seemed like Konilo had a pretty good self-contained ecosystem in place already as an editor. The interactive aspect of a Forth also seems in line with the aspect of composing music.
\par

So, developing a serial command language is going to be more important than hooking things up to Konilo. It would make sense to focus on this more, than actually hooking it up to Konilo. Actually, it's a prerequisite. By the time things are ready to actual integrate with Konilo, the serial command language should already be in place.
\par

Working on this task should mostly be about working with the ergonomics of the Konilo ecosystem and making it work with Sonilo command language.
\par

I am developing a new task now.
\par

Well, I've made a new task, but I think I need code in functionality to re-import tasks into a grouped schedule.
\par

{\medskip\noindent\bf\font\titlefont=cmitt12 at 17pt\titlefont 2025-11-11}\medskip
{\smallskip\noindent\bf\font\titlefont=cmr12 at 13pt\titlefont $\bullet$ 16:50: sonilo: serial protocol concepts, penciling}\smallskip
{\smallskip\noindent\bf\font\titlefont=cmr12 at 13pt\titlefont $\bullet$ 20:25: sonilo: serial protocol concepts, inking}\smallskip
{\medskip\noindent\bf\font\titlefont=cmitt12 at 17pt\titlefont 2025-11-13}\medskip
{\smallskip\noindent\bf\font\titlefont=cmr12 at 13pt\titlefont $\bullet$ 10:37: Konilo feels very far off tbh. best I can do is IO experiments}\smallskip
There's just so much I don't know about the konilo ecosystem. how does it load files? can I use source control? is it the right langauge abstraction? can it handle all the ugens I want to use?
\par

It looks like ultimately I'll be building programs to generate basic msgpack structures. So, maybe that can be a focus.
\par

{\medskip\noindent\bf\font\titlefont=cmitt12 at 17pt\titlefont 2026-01-02}\medskip
{\smallskip\noindent\bf\font\titlefont=cmr12 at 13pt\titlefont $\bullet$ 15:43: check-in}\smallskip
Ilo probably won't be used for constructing ugen chains, due to the memory overhead. But, Konilo is going to be important for building abstractions. There's not much more to say that hasn't been said already. Ilo has the ability to do I/O commands, and getting it to work with sonilo would mean making a virtual serial device, and then calling those commands from Konilo.
\par

\xrdef{node6}
{\medskip\noindent\bf\font\titlefont=cmss17 at 20pt\titlefont mem}\medskip
{\medskip\noindent\bf\font\titlefont=cmitt12 at 17pt\titlefont 2025-11-13}\medskip
{\smallskip\noindent\bf\font\titlefont=cmr12 at 13pt\titlefont $\bullet$ 10:25: memory allocation: core components close to done, how to tie together?}\smallskip
The buddy system is close to being done. The block stack API is finished. The cache is is finished. These are components that need to work together to build a cohesive memory allocation system.
\par

For the small memory allocation system to work, the buddy system needs to automatically request new blocks under the hood.
\par

The block-level allocation system will be used for allocating output blocks.
\par

More thoughts needed here.
\par

{\medskip\noindent\bf\font\titlefont=cmitt12 at 17pt\titlefont 2025-11-28}\medskip
{\smallskip\noindent\bf\font\titlefont=cmr12 at 13pt\titlefont $\bullet$ 09:59: memory primitives: alloc, block, cleanup}\smallskip
alloc: allocate N words (N < 64) block: allocate a block
\par

cleanup: frees all memory in the current context
\par

\xrdef{node7}
{\medskip\noindent\bf\font\titlefont=cmss17 at 20pt\titlefont sonilo command language}\medskip
{\medskip\noindent\bf\font\titlefont=cmitt12 at 17pt\titlefont 2025-11-13}\medskip
{\smallskip\noindent\bf\font\titlefont=cmr12 at 13pt\titlefont $\bullet$ 10:33: msgpack, what is the initial command set?}\smallskip
I'm decided to use msgpack as a communicaiton protocol. Or a subset of msgpack at least. Using something that has a spec already written out has advantages. Having data types like arrays and objects could come in handy down the line.
\par

The truth is, I don't know exactly where this is headed, and I need to design with some flexibility in mind.
\par

msgpack is also handy because I can write intermediate programs to read (and possibly generate, provided it fits the subset of msgpack I'll support) msgpack programs.
\par

msgpack programs might be the low-level layer for generating the first sonilo sounds.
\par

{\medskip\noindent\bf\font\titlefont=cmitt12 at 17pt\titlefont 2026-01-02}\medskip
{\smallskip\noindent\bf\font\titlefont=cmr12 at 13pt\titlefont $\bullet$ 14:50: check-in}\smallskip
More recently, I've found this to be a more important task. This is going to be the way I get sonilo into other systems like Konilo.
\par

I've been wondering if maybe I should just get my silly test language working as a protocol, just to get a feel for the system and working with Konilo.
\par

I'm still somewhat optimistic about working with Konilo. I think I may want to build an SDL backend for Konilo at some point to allow for primitive drawing (1-bit stuff only).
\par

{\medskip\noindent\bf\font\titlefont=cmitt12 at 17pt\titlefont 2026-01-16}\medskip
{\smallskip\noindent\bf\font\titlefont=cmr12 at 13pt\titlefont $\bullet$ 09:26: test suite as prototype for command language?}\smallskip
More and more, it's looking like this test suite language is acting like a rough draft for what the command language will look like.
\par

I think the one 32-bit register idea is quirky, but it acts as a good focus point for design. I think I'd probably built out more general purpose operations that work on that register, doing things equivalent to stack operations you'd find in a forth.
\par

{\medskip\noindent\bf\font\titlefont=cmitt12 at 17pt\titlefont 2026-01-31}\medskip
{\smallskip\noindent\bf\font\titlefont=cmr12 at 13pt\titlefont $\bullet$ 16:01: task ordering updates, ugens child of some other tasks now}\smallskip
ilo interpreter stuff needs to at least happen after ugen stuff. And it's a prereq for any gesture stuff.
\par

By the time unit generators are finished, I think thinking about using a higher-level abstraction like Konilo may be warranted. At that point, it'll start to feel possible to build demos.
\par

\xrdef{node8}
{\medskip\noindent\bf\font\titlefont=cmss17 at 20pt\titlefont phasor}\medskip
{\medskip\noindent\bf\font\titlefont=cmitt12 at 17pt\titlefont 2025-11-13}\medskip
{\smallskip\noindent\bf\font\titlefont=cmr12 at 13pt\titlefont $\bullet$ 10:39: technically this is today's task, ha. I am so behind}\smallskip
The DSP algorithm for this is already available in sonilo. The work involved here will reworking the memory so it reads from sonilo's fixed memory system.
\par

The DSP data will have single entry point, a memory address represented by a virtual pointer. In this address will be data containing pointers to inputs (constants or blocks), and outputs (blocks). There will also maybe a word or two for storing state the current phase.
\par

I can probably write most of this without needed the actual DSP memory fully done.
\par

{\medskip\noindent\bf\font\titlefont=cmitt12 at 17pt\titlefont 2025-11-24}\medskip
{\smallskip\noindent\bf\font\titlefont=cmr12 at 13pt\titlefont $\bullet$ 10:13: let's try porting phasor today}\smallskip
I am so behind on the fundamental memory IO and bytecode interpreter stuff (AKA the "hard bits"). Phasor is mostly just moving pre-existing code around and setting up the scaffolding.
\par

Got some initial phasor DSP code added. untested, but it's mostly a copy-paste port.
\par

A phasor exists in a 4 word block. The first 2 words are the ports. It has one input port for frequency, and one output port for the returned signal.
\par

I don't have a great way to hook this up to anything at the moment. It would be trivial to write a small C program for this, but that's not too useful. Building towards ILO is going to be more useful because it will allow me to write small test programs that compose different unit generators together.
\par

Creating a function `ugen\_phasor`, which should emulate the behavior of the phasor unit generator, popping a parameter off the stack to set the input port, and then pushing a parameter back onto the stack.
\par

The abstractions don't feel quite there yet, but I think it's good enough. I really do need to working on getting the ilo interpeter working now.
\par

I'm going to say this is "done", just for the sake of progress.
\par

{\smallskip\noindent\bf\font\titlefont=cmr12 at 13pt\titlefont $\bullet$ 12:39: Small C program it is}\smallskip
I've spent the last hour brainstorming the ILO code. The problem is more complex than I thought, requiring me to implement many moving parts. The DSP code does not need that yet. It's going to be way faster just to write the silly C program. Later, it can be rewritten to be an ilo ROM.
\par

{\smallskip\noindent\bf\font\titlefont=cmr12 at 13pt\titlefont $\bullet$ 12:49: writing a silly C program}\smallskip
I'm so behind on things, I just gotta confirm these small bits work.
\par

Wow, this blows up again in complexity. My basic stack struct has bugs? ugh.
\par

Ah, I see. Phasor needed to know what the output pointer was going to be. Initialization functions need to take in variables.
\par

That's all the time we have, unfortunately. 
\par

{\medskip\noindent\bf\font\titlefont=cmitt12 at 17pt\titlefont 2025-11-25}\medskip
{\smallskip\noindent\bf\font\titlefont=cmr12 at 13pt\titlefont $\bullet$ 09:21: more work with the bits surrounding phasor}\smallskip
DSP code took about 5 minutes, designing a flexible interface took up the remaining time.
\par

Converting the bytes in floats to ints is proving to take some time. I really just need a helper function for this.
\par

we got a working phasor written to disk finally
\par

{\smallskip\noindent\bf\font\titlefont=cmr12 at 13pt\titlefont $\bullet$ 10:03: cleanup: build out APIs, phasor init/compute should use stack interface}\smallskip
arrays (for stack), and parameters. oh, also float/int conversion
\par

I'm not going to try and over-categorize. I'm throwing everything into a "util" file for the being, and this will live in the DSP folder because only DSP utilities are going to be using this stuff (at least in the short term).
\par

{\smallskip\noindent\bf\font\titlefont=cmr12 at 13pt\titlefont $\bullet$ 10:38: phasor, officially done}\smallskip
\xrdef{node9}
{\medskip\noindent\bf\font\titlefont=cmss17 at 20pt\titlefont initial dsp}\medskip
{\medskip\noindent\bf\font\titlefont=cmitt12 at 17pt\titlefont 2025-11-24}\medskip
{\smallskip\noindent\bf\font\titlefont=cmr12 at 13pt\titlefont $\bullet$ 11:27: writing out pseudo-code for an ILO program computing phasor}\smallskip
This exercise is exposing some points of friction for me.
\par

The hard part is memory structure, once again. The ILO bits are very straight forward, once the memory structure is set up.
\par

Point of friction: I started this task wanting to just write out pseudo-assembler code for ILO. Then I realized, I need memory addresses. How do I generate those things? I'd need to allocate and initialize data for the phasor struct, in addition to the parameter.
\par

So, to make a program is really a collaboration between the memory writer and the ilo assembler. I'd use the memory writer to set up the memory structure, and the assembler to generate ilo bytecode.
\par

Ideally, it would be nice for programs to use the original ilo C implementation to start. I can make it so the ilo program is at the start of the file, with all the special memory after it. The program would terminate by a goto that goes beyond the ilo bounds.
\par

ilo programs are going to be very small, and the address space localized to just a handful of blocks. This phasor program would exist on one block.
\par

What does this phasor program do? It opens a file for writing, runs a loop writing stuff to that file, then closes that file. In the loop, the phasor ugen is called to generate a block of signal. That signal then gets written to disk. It is assumed that the phasor data is written to disk pre-initialized.
\par

Okay, this is clearly more work than I expected.
\par

\xrdef{node10}
{\medskip\noindent\bf\font\titlefont=cmss17 at 20pt\titlefont blsaw}\medskip
{\medskip\noindent\bf\font\titlefont=cmitt12 at 17pt\titlefont 2025-11-29}\medskip
{\smallskip\noindent\bf\font\titlefont=cmr12 at 13pt\titlefont $\bullet$ 08:39: blsaw initial code}\smallskip
Just porting the initial code and getting it in using some of the conventions established in phasor.
\par

\xrdef{node11}
{\medskip\noindent\bf\font\titlefont=cmss17 at 20pt\titlefont ugen bytecode}\medskip
{\medskip\noindent\bf\font\titlefont=cmitt12 at 17pt\titlefont 2025-11-30}\medskip
{\smallskip\noindent\bf\font\titlefont=cmr12 at 13pt\titlefont $\bullet$ 08:59: ilo interpreter can wait, use custom bytecode format instead}\smallskip
I think the DSP is specialized enough that I can get away with it.
\par

Ilo, on the other hand, is going to be more important in situtaions where small, inline executable programs are desirable. I'm mainly thinking this could be used for  gesture iterators.
\par

{\smallskip\noindent\bf\font\titlefont=cmr12 at 13pt\titlefont $\bullet$ 09:50: New task created}\smallskip
Instead of trying to use ILO, I'm going to make my own format. It will be more compact, simpler to build, and less IO overhead.
\par

The proposed format is also more data than executable code, which is good at this stage.
\par

{\medskip\noindent\bf\font\titlefont=cmitt12 at 17pt\titlefont 2026-01-18}\medskip
{\smallskip\noindent\bf\font\titlefont=cmr12 at 13pt\titlefont $\bullet$ 10:10: review: I might be able to start this}\smallskip
There's a buddy allocator that works on a block of memory that works as far as I can tell, and that might be enough to jump into things.
\par

Uh oh, where did I write this down.
\par

Oh it was in the ilo interpreter task.
\par

I was more vague about this than I'd like. A similar packing format to Ilo, except that there'd be 4-bit opcodes instead of 8-bit opcodes, packing up to 8 instructions in a single word. Similar to Ilom, the instructions would be able to move the program counter to read words.
\par

I think I was initially thinking there'd be stack operations to push/pop parameters. But, now I'm wondering if I just want to just emulate the patchwerk/graforge setup I had. In that setup, the structure was a list of nodes, with all the memory locations baked in.
\par

What does a command language look like for a list of nodes?
\par

A node has data, which is stored in an address, and it has some kind of opcode which processes that data. Memory is a 16-bit address, so an opcode lookup id can be 16 bits to make both fit on a single word. A-Z fits in 5 bits with 6 extra spaces to spare, so 16 bits gives 3 letters with an extra bit. alphanumeric would be nice. but 6-bit gives us 2 letters. some kind of custom base 36 radix format could possibly yield 3. maybe base 26 could yield 4 letters. But, if the opcode IDs were arbitrary, that's still a lot of unit generators.
\par

Blocks might still be present here (maybe they get abstracted away?) But if they aren't, jump instructions will be needed to leap from block to block. Gosh, that sounds messy. I think I'm going to want some kind of virtual array for this. I already came up with a scheme for this way back in september.
\par

This could get reduced even further: no VM, just data. One word is a node. A block contains up to 64 nodes. A command could be written to execute the word register containing the opcode id and the data. It could write stuff to memory and the word register.
\par

The nice thing about having a command like that is that instruction fetching could be "upgraded" to read from a different data structure other than a raw block.
\par

Running a single instruction is one thing, but what about a block of instructions? What if I want to run that block many many times (as is the case with DSP blocks)? I need loops, and this test program isn't the right place for that. Konilo/Sonilo could make some sense too. But I think I'll end up coding something up in C.
\par

the data/opcode paradigm feels pretty solid. But, I'll need a way to build up that data. That becomes thew new challenge. But certainly it would be trivial to make a simple opcode that reads in data and adds 1 to it or something. That would be enough to get a general system in place.
\par

{\medskip\noindent\bf\font\titlefont=cmitt12 at 17pt\titlefont 2026-01-24}\medskip
{\smallskip\noindent\bf\font\titlefont=cmr12 at 13pt\titlefont $\bullet$ 13:48: ugen bytecode initial scoping}\smallskip
Thinking about subtasks here
\par

core primitive is an instruction that fits into a 32-bit word. 16 bits for data, 16 bits for the opcode.
\par

Things like ugen structure, offline rendering, etc, are already well defined tasks that are parents of this task.
\par

When the single instruction primitive is created, bulk instruction execution needs to happen. I'm thinking about executing an array of instructions. That data structure is already built.
\par

\xrdef{node12}
{\medskip\noindent\bf\font\titlefont=cmss17 at 20pt\titlefont allocator}\medskip
{\medskip\noindent\bf\font\titlefont=cmitt12 at 17pt\titlefont 2025-11-30}\medskip
{\smallskip\noindent\bf\font\titlefont=cmr12 at 13pt\titlefont $\bullet$ 10:01: make task for memory allocator}\smallskip
This will be a high-level allocator, wrapping the block list, bitset, and the buddy allocator.
\par

got caught up with printing it to a PDF
\par

{\medskip\noindent\bf\font\titlefont=cmitt12 at 17pt\titlefont 2026-01-02}\medskip
{\smallskip\noindent\bf\font\titlefont=cmr12 at 13pt\titlefont $\bullet$ 15:33: check-in}\smallskip
since writing this, I've gotten a more formed idea of what an "allocator" is going to be. I will articulate what I know now.
\par

Memory management is a layered ordeal. The memory to be managed is $2^16$ words of 32-bit. Memory is broken up into chunks of 64-words, called blocks. Free blocks are maintained on a stack. Allocation is a matter of popping the block, and pushing the block back is freeing the memory.
\par

Entire blocks are typically used as audio buffers to store signals. Hence the 64 size, which is a pretty common block size for audio systems.
\par

A special allocator is used to allocate small segments of memory from the buddy system. This is a buddy allocator. Typically, this is used to allocate memory for unit generators.
\par

Multiple buddy allocators are used together to form the main allocator used for small allocations.
\par

When a block is retrieved from the block stack, it is stored in a set of blocks. This makes it easier to free memory later, because it frees everything at once.
\par

\xrdef{node13}
{\medskip\noindent\bf\font\titlefont=cmss17 at 20pt\titlefont docgen}\medskip
{\medskip\noindent\bf\font\titlefont=cmitt12 at 17pt\titlefont 2025-11-30}\medskip
{\smallskip\noindent\bf\font\titlefont=cmr12 at 13pt\titlefont $\bullet$ 10:57: creating a docgen for my sonilo input.txt file}\smallskip
I need a better way to read a review this stuff. attempting to generate tex for this
\par

\xrdef{node14}
{\medskip\noindent\bf\font\titlefont=cmss17 at 20pt\titlefont buddy slots}\medskip
{\medskip\noindent\bf\font\titlefont=cmitt12 at 17pt\titlefont 2025-12-03}\medskip
{\smallskip\noindent\bf\font\titlefont=cmr12 at 13pt\titlefont $\bullet$ 21:49: initial buddy slots work}\smallskip
{\medskip\noindent\bf\font\titlefont=cmitt12 at 17pt\titlefont 2025-12-12}\medskip
{\smallskip\noindent\bf\font\titlefont=cmr12 at 13pt\titlefont $\bullet$ 21:18: checking in buddy slots}\smallskip
Tonight I feel like my abstractions are sloppy. I'm building up to a more sophisticated allocator, but things like the block list and the parameter stack need to be more out of the way.
\par

No real work done. I feel stuck.
\par

{\medskip\noindent\bf\font\titlefont=cmitt12 at 17pt\titlefont 2025-12-13}\medskip
{\smallskip\noindent\bf\font\titlefont=cmr12 at 13pt\titlefont $\bullet$ 10:11: morning thoughts: missing abstractions}\smallskip
This morning, I still feel that the "multi-buddy" allocator is missing abstractions. The point of friction is with my current system there is no way to build up abstractions. Something like this allocator needs to be able to have a pointer to one "thing", which in turn is a collection of pointers to a few other "things". To do this, you need an allocator, but this is what is built, and thus there is a circular dependency.
\par

So, memory needs to be allocated for the allocator. This is a bootstrapping problem. So, I'm thinking there may be an incoming intermediate abstraction required between the global block list and this multi-buddy allocator.
\par

My test system uses a single "cursor" as a pointer in memory. This has an issue with abstractions. Say I have a cursor pointing to a compound data structure containing a list of pointers to many things, with one pointing to a block of memory. I may have an operation that needs to set that cursor to that block of memory. Once I do the operation, I need a way to jump back to that compound data structure.
\par

A design I think I remember reading about in the Elements of Computing Systems is the idea of having two values that can be swapped, with only the active one being able to be set.
\par

So, for this hypothetical block operation, the commands would look like: set word register to be desired block address, swap cursors, set cursor from word, do block operation, swap back.
\par

Bootstrapping would involve setting up the block list, and then initializing a zero page with the address of the block list as the first item on the zero page. From there, you could set up a blockset and a stack to be addresses on the zero page. So on and so forth.
\par

{\smallskip\noindent\bf\font\titlefont=cmr12 at 13pt\titlefont $\bullet$ 21:15: rework integer parsing to insert at LSB}\smallskip
Emerged from my work on zero pages. Basically, I found a clever way to write words to specific slots in a block without the need for a stack (a stack would have been necessary for an explicit write word to block operation because it'd need at least two args). The approach requires storing input from the LSB instead of from the MSB.
\par

huzzah done
\par

{\medskip\noindent\bf\font\titlefont=cmitt12 at 17pt\titlefont 2025-12-16}\medskip
{\smallskip\noindent\bf\font\titlefont=cmr12 at 13pt\titlefont $\bullet$ 21:02: before buddy slots, a zero page}\smallskip
A buddy slots system is going to need a stack to handle multi-variable function calls. I need to get enough of a system working so that there is a stack to use.
\par

{\smallskip\noindent\bf\font\titlefont=cmr12 at 13pt\titlefont $\bullet$ 21:29: Ugh, what are the arguments need to initialize this multi-buddy system?}\smallskip
It needs to be able to allocate new blocks when needed. So, that requires the block stack, probably via the bitset.
\par

{\medskip\noindent\bf\font\titlefont=cmitt12 at 17pt\titlefont 2026-01-02}\medskip
{\smallskip\noindent\bf\font\titlefont=cmr12 at 13pt\titlefont $\bullet$ 15:31: buddy-slots check-in}\smallskip
still blocked. this is a more complex function made up smaller parts (a sonilo first), so some more complex components are needed (a zero page).
\par

{\medskip\noindent\bf\font\titlefont=cmitt12 at 17pt\titlefont 2026-01-15}\medskip
{\smallskip\noindent\bf\font\titlefont=cmr12 at 13pt\titlefont $\bullet$ 14:37: what precisely do I need to initialize the buddy slot allocator?}\smallskip
It needs to allocate blocks when needed and put them in the bitset. That's one 32-bit word as it turns out.
\par

That might be it for initialization.
\par

If I'm doing the math correctly: a block can house 6 instances of a buddy allocator (10 words each), 3 words are needed to construct the args needed to call a single buddy allocator. 1 word is left to store the address of the bitset/blockstack combo.
\par

0-59: buddy slots. 60: address of block stack. 61-63: space to dynamically fill out buddy slots.
\par

{\smallskip\noindent\bf\font\titlefont=cmr12 at 13pt\titlefont $\bullet$ 19:32: brainstorm: buddyslots function composition (part 1)}\smallskip
{\smallskip\noindent\bf\font\titlefont=cmr12 at 13pt\titlefont $\bullet$ 21:22: brainstorm: buddyslots function composition}\smallskip
The main idea here was to imagine a interface for buddy slots that can somehow use the existing functions in my test suite for allocation (it's looking like my sonilo command language will closely resemble this).
\par

I'll have to look at what goes into freeing things, as that may require tracking the k size.
\par

Once again, there's the matter of address space mapping. This time, it's mapping the addresses returned by the buddy allocator (relative word indices 0-63) to words in memory. With the buddy slot system, there'd need to be a way to differentiate between which instance of the buddy system the index belongs to.
\par

{\medskip\noindent\bf\font\titlefont=cmitt12 at 17pt\titlefont 2026-01-16}\medskip
{\smallskip\noindent\bf\font\titlefont=cmr12 at 13pt\titlefont $\bullet$ 10:57: re-examining buddy system API for buddy slots design}\smallskip
The buddy system uses a stack that is instantiated, which is actually good news. It should be a straightforward thing to create new operations that work with the existing operations.
\par

For example, say the buddy system data is stored at address {\tt 0x123}. To allocate 4 words, $k = 2$, and the memwrite code would look like {\tt \#2 aa \#123 aa ma}. One could build an operation called {\tt ss} that selects a slot and generates data for the buddy allocator to read, which then can be pushed onto the stack. So, using the allocator with the buddy system could look something like: {\tt \#2 aa sc \#0 ss sc aa ma}. Here, the two cursors are used to jump between the stack and the buddy slot data.
\par

The alt cursor might need to be reserved for the zero page, and the main cursor would be used for the stack. In such a situation, perhaps the address of the buddy slot is stored in the word slot 2 of the zero page. Swap to the zero page, read the slot (jump forward 2, read, jump backward 2), swap to the stack, push the value onto the stack. It could look like this: {\tt \#2 sc \#0 aa \#2 jf rd \#2 jb sc aa ss aa ma}. In this situation, {\tt ss} uses the stack entirely for both the address of the buddy slots and which slot to select.
\par

The stack usage of latter approach with {\tt ss} is more consistent with with the design {\tt ma}, so I'm going to go with that.
\par

{\medskip\noindent\bf\font\titlefont=cmitt12 at 17pt\titlefont 2026-01-17}\medskip
{\smallskip\noindent\bf\font\titlefont=cmr12 at 13pt\titlefont $\bullet$ 20:36: buddy slot system: append and selection operations}\smallskip
The bottoms-up thing would be to implement a 'buddy slot select' operation, {\tt ss}. In order to select a slot, a the slots need to be initialized {\tt sm} (slots make) and a slot needs to be appended to {\tt sx} (slot extend).
\par

These should all use a stack, and the stack should be allocated by the zero page.
\par

I got as far as setting up a stack. exhausting work. I forgot about the trick of not including '\#' for jf/jb commands.
\par

\xrdef{node15}
{\medskip\noindent\bf\font\titlefont=cmss17 at 20pt\titlefont zero page}\medskip
{\medskip\noindent\bf\font\titlefont=cmitt12 at 17pt\titlefont 2025-12-16}\medskip
{\smallskip\noindent\bf\font\titlefont=cmr12 at 13pt\titlefont $\bullet$ 21:10: setting up zero page task}\smallskip
{\smallskip\noindent\bf\font\titlefont=cmr12 at 13pt\titlefont $\bullet$ 21:16: writing initial words on zero page}\smallskip
I'm going to roughly transcribe some of what I was thinking a few days ago with this zero page.
\par

A zero page is like some kind of boot sector. The idea is that there's this top-level block that can be used to read and write data.
\par

The thinking is this can bootstrap an initial system for memory allocation. So that means it stores a pointer to the global block stack, sets of blocks it has allocated, a stack, amongst other possible things.
\par

What's of particular importance is getting a stack working. When that happens, I can build up commands that take more than one argument. This is needed for my buddy system, amongst other things.
\par

So this bootstrapping process. How does it work? Set up the block stack. Allocate a zero page from the block stack. Set up a bitset for memory, place in block 1. Allocate a stack (make sure to save it in the bitset), place in block 2.
\par

I don't know if the bitset and the block stack are working together yet.
\par

These are 16-bit addresses in a 32-bit word space. Might be worth putting the bitset and the block stack together. That way, I can have a function that in one go pops a block and stores it in a bitset.
\par

The first thing that needs to be made is something that initializes a zero page, "zp". ZP will take in the address of a block stack, use that to allocate a new block, zero it out, and set the first words to be the address of the block stack. That roughly would look like "\#b go bi zp go" where "\#b" is the address of the block stack (address 0?).
\par

After that, the process of setting up bitset, and storing that in the zero page.
\par

The cursor is currently set to the zero page, and it it needs to read the address of the block stack from slot 0. To do this, a "switch cursor" functionality is needed, "sc". This allows the system to jump between two cursors. Once the block stack is selected, it can allocate a new block and use that to be a set to store blocks. The commands to do this are "rd sc go ba go si".
\par

With the set initialized it must be stored in slot 1 of  the zero page. To do this, the memory address for slot 1 must be computed. We assume that blocks are memory aligned, so the start of the block points to slot 0, the next block points to slot 1, etc. So, go to the zero page, move forward 1. write the data, jump back (different from what I was thinking about a few days ago). "cf" and "cb" for cursor forward and cursor back aren't taken, so we'll use these. Code I believe works out to be "sc \$1 cf sc rc sc wr \$1 cb".
\par

Get the stack initialized, write to slot 2. "rd sc go ba ai go sc \$2 cf rc sc wr \$2 cb"
\par

I'm thinking by necessity, the set will need to be written into the unused MSB half of slot 0, and to create another function, say "bb" that allocates a block and stores that address in the register. Maybe "aa" is some kind of function that merges the address in the cursor and the address in the word register together. After allocating the bitset and jumping to that address, but before switching back to the zero page: "rc sc aa wr" will copy the set addr to the word, switch to the zero page address containing the block stack, merge the addresses, and then write back to that top slot. Allocating the stack would use "bb" instead of "ba", and store in slot 1 instead of 2: "rd sc go bb ai go sc \$1 cf rc sc wr \$1 cb"
\par

That's enough rambling for today.
\par

{\medskip\noindent\bf\font\titlefont=cmitt12 at 17pt\titlefont 2025-12-17}\medskip
{\smallskip\noindent\bf\font\titlefont=cmr12 at 13pt\titlefont $\bullet$ 08:24: set up some initial zero page stuff}\smallskip
New commands: {\tt zp} for "zero page", {\tt cf} and {\tt cb} (alt: {\tt jf} and {\tt jb}) for "cursor jump forward/backward", {\tt sc} for "switch cursor", {\tt aa} for "merge two 16-bit addresses into one word", {\tt bb} for "for 32-bit word (MSB, LSB), allocate a block off the stack in address in LSB, then mark the address in the bitset in MSB".
\par

Writing up some initial TODOs in my test suite.
\par

Some potential tests written out in my input.txt file. Another command I'm going to add: "sn".
\par

This is why things take so long. Tasks here cascade into smaller tasks that individually don't warrant building a subtask for (currently too much overhead required to build a task). My TDD mindset is slowing things down too. But I think it'll be worth it when I port this to rust.
\par

{\medskip\noindent\bf\font\titlefont=cmitt12 at 17pt\titlefont 2025-12-18}\medskip
{\smallskip\noindent\bf\font\titlefont=cmr12 at 13pt\titlefont $\bullet$ 09:55: implement zero page command 'zp'. blocklist blockers.}\smallskip
the zero page has been implemented, but there are some blockers related to the blocklist. Words to follow
\par

This block list addressing scheme is going to need to get looked at again.
\par

Currently I've been assuming that the block list is initialized at address 0. Addresses are stored as memory aligned block addresses. When I run 'gb', it assumes that block '0' is the first free block available. so, it skips the words used to make the blocklist. However, my test doesn't have that assumption.
\par

I think it'll work out okay for today, but yeah, may want to think again about what this block list is doing.
\par

oof. I think there's weird overflow bugs happening. I think my first address to be returned is block 0x400, the last block in memory. it works backwards. When I add a bias to that, it overflows. The issue is the stack assumes all 1024 blocks are available, but by definition, they are not because some of those blocks (16 of them) are used for this block list.
\par

Okay, so if a block list can be initialized to any address, what exactly are those 10-bit numbers mapping to? You can't just turn that into a memory address by multiplying by 64 and adding a 1024 bias, that only works when the memory is 0.
\par

So, the so-called "block addresses" need to be thought of as a mapping from a 10-bit number to a 16-bit word address.
\par

Those 10-bit numbers are offsets local to the address of the block stack. Brute force, the mapping could be a lookup table. 1024 16 bit numbers packed into 32 words means 512 words, which would be 8 blocks. Perhaps there's a better way.
\par

If we assume the mapping is mostly linear, but skip the 16 blocks used for the block stack, that mapping could be realized as a closed-form expression.
\par

It'll never be 1024 blocks, only 1008. That's 15 blocks exactly, not 16. Might be worth thinking about.
\par

making a task for this.
\par

{\medskip\noindent\bf\font\titlefont=cmitt12 at 17pt\titlefont 2026-01-02}\medskip
{\smallskip\noindent\bf\font\titlefont=cmr12 at 13pt\titlefont $\bullet$ 15:26: check-in: zero page, still blocked by blocklist}\smallskip
Apparently, my zero page unit tests are breaking because the blocklist is broken, so that component is getting rewritten. Not much more to say right now.
\par

{\medskip\noindent\bf\font\titlefont=cmitt12 at 17pt\titlefont 2026-01-15}\medskip
{\smallskip\noindent\bf\font\titlefont=cmr12 at 13pt\titlefont $\bullet$ 10:35: zero page test fixes}\smallskip
trying to remember where I left off
\par

distracted by docgen.
\par

almost second guessing my block-to-memory calculations. If my blocklist starts at address 0x10, the first block is going to be at 0x180 because 0x10 is in the first five blocks. You're If the blocklist was stored in 0x11, it would also be 180. I think it's block-level alignment (for the most part).
\par

not going to overthink this. It's good enough. moving on.
\par

{\smallskip\noindent\bf\font\titlefont=cmr12 at 13pt\titlefont $\bullet$ 13:32: what's next? ramblings.}\smallskip
I've just gotten the zero page set up with the block stack.
\par

According to my notes, bitset + block stack should probably happen next.
\par

I used to have a "bw" that would convert a BBB.II word into a word address based on block and offset. I don't think that can work anymore.
\par

The bottom line is: can the zero page still be used to generate the complex data structure needed for the buddy slot?
\par

It doesn't seem like I have a concrete answer for this written down anywhere.
\par

{\smallskip\noindent\bf\font\titlefont=cmr12 at 13pt\titlefont $\bullet$ 15:43: implementing jumps (jf and jb)}\smallskip
The thinking I have here is that 'jf' and 'jb' will be used to do local jumps forward/backward relative to the current cursor position. It'll treat the word register like a stack and read the last 2 nibbles as an offset.
\par

{\smallskip\noindent\bf\font\titlefont=cmr12 at 13pt\titlefont $\bullet$ 16:26: working on 'sc'}\smallskip
I have run into a weird printing issue. troubleshooting.
\par

OH! forgot to start the comment. It was just parsing all those characters like a program I think.
\par

{\smallskip\noindent\bf\font\titlefont=cmr12 at 13pt\titlefont $\bullet$ 16:51: pack addresses ('pa', formerly 'aa')}\smallskip
This sort of thing is needed to merge the bitset and the global blocklist addresses into one word. I was going to use 'aa', but that's already taken (array append).
\par

{\smallskip\noindent\bf\font\titlefont=cmr12 at 13pt\titlefont $\bullet$ 17:05: 'sn' get number of elems for set}\smallskip
I forget why I needed this one. Probably testing? I'll do it anyways.
\par

{\smallskip\noindent\bf\font\titlefont=cmr12 at 13pt\titlefont $\bullet$ 17:19: 'bb' allocate and mark block}\smallskip
Note to self: block address space is 1-1023 inclusive. 0 is not used. the bitset can include 0, but that's being reserved to make the linked list logic run smoother.
\par

{\smallskip\noindent\bf\font\titlefont=cmr12 at 13pt\titlefont $\bullet$ 17:45: I think I have it all done. this zero page work is complete.}\smallskip
\xrdef{node16}
{\medskip\noindent\bf\font\titlefont=cmss17 at 20pt\titlefont blockstack improvements}\medskip
{\medskip\noindent\bf\font\titlefont=cmitt12 at 17pt\titlefont 2025-12-18}\medskip
{\smallskip\noindent\bf\font\titlefont=cmr12 at 13pt\titlefont $\bullet$ 09:55: implement zero page command 'zp'. blocklist blockers.}\smallskip
the zero page has been implemented, but there are some blockers related to the blocklist. Words to follow
\par

This block list addressing scheme is going to need to get looked at again.
\par

Currently I've been assuming that the block list is initialized at address 0. Addresses are stored as memory aligned block addresses. When I run 'gb', it assumes that block '0' is the first free block available. so, it skips the words used to make the blocklist. However, my test doesn't have that assumption.
\par

I think it'll work out okay for today, but yeah, may want to think again about what this block list is doing.
\par

oof. I think there's weird overflow bugs happening. I think my first address to be returned is block 0x400, the last block in memory. it works backwards. When I add a bias to that, it overflows. The issue is the stack assumes all 1024 blocks are available, but by definition, they are not because some of those blocks (16 of them) are used for this block list.
\par

Okay, so if a block list can be initialized to any address, what exactly are those 10-bit numbers mapping to? You can't just turn that into a memory address by multiplying by 64 and adding a 1024 bias, that only works when the memory is 0.
\par

So, the so-called "block addresses" need to be thought of as a mapping from a 10-bit number to a 16-bit word address.
\par

Those 10-bit numbers are offsets local to the address of the block stack. Brute force, the mapping could be a lookup table. 1024 16 bit numbers packed into 32 words means 512 words, which would be 8 blocks. Perhaps there's a better way.
\par

If we assume the mapping is mostly linear, but skip the 16 blocks used for the block stack, that mapping could be realized as a closed-form expression.
\par

It'll never be 1024 blocks, only 1008. That's 15 blocks exactly, not 16. Might be worth thinking about.
\par

making a task for this.
\par

{\medskip\noindent\bf\font\titlefont=cmitt12 at 17pt\titlefont 2025-12-19}\medskip
{\smallskip\noindent\bf\font\titlefont=cmr12 at 13pt\titlefont $\bullet$ 10:00: block stack v2 thoughts}\smallskip
{\medskip\noindent\bf\font\titlefont=cmitt12 at 17pt\titlefont 2026-01-02}\medskip
{\smallskip\noindent\bf\font\titlefont=cmr12 at 13pt\titlefont $\bullet$ 14:10: check-in}\smallskip
I'm not sure where I'm at here. My logs are blank. I think I added some tests.
\par

hang on, docgen is busted... 
\par

okay looking at the zero page task. lots of rambling there
\par

I think this is what the blockstack improvements are: changing the data to use packed 10-bit numbers, and creating a way to turn those numbers to memory addresses.
\par

the problem of packing 10-bit ints into 32-bit word memory is tricky due to alignment. I'm wondering if it might be beneficial to generalize this problem. It turns out that the entire memory can be addressable as a bit array. Sonilo memory has $2^16$ 32-bit words, which means that there are $2^16 2^5 = 2^21$ addressable bits, which fits inside of a 32-bit number.
\par

There needs to be a way to convert blockstack addresses to word addresses. It's a kind of linear mapping, but it skips the range of addresses needed for the block stack.
\par

The 'ba' command returns block addresses, which are just those 10 bit numbers. Another utility factoring, in the block stack address should be used for conversion. Not sure what that looks like though.
\par

{\medskip\noindent\bf\font\titlefont=cmitt12 at 17pt\titlefont 2026-01-04}\medskip
{\smallskip\noindent\bf\font\titlefont=cmr12 at 13pt\titlefont $\bullet$ 20:23: global bit array brainstorm}\smallskip
first thing I'm going to do is change the blockstack to use 10-bit ints. to do this, I am going to generalize the problem. packing the ints becomes a lot easier when the underlying data structure is a bit array and the word boundary logic is abstracted away.
\par

I would like to implement a global bit array. Every bit can be individually addressed in this memory space using 32-bit integers.
\par

It's an array, so the core operations are some kind of "read" and some kind of "write", to some location.
\par

Usually, it'll be the case that multiple bits will want to be read or written to at once, such as is the case for these 10-bit ints. 32-bits would seem to be the upper limit for this.
\par

so, for read and write, a start (offset) position and size would be the arguments. the data to be written to a 32-bit integer.
\par

ah, I've already written about this. I think I might transcribe my outline, and see how long this takes.
\par

{\smallskip\noindent\bf\font\titlefont=cmr12 at 13pt\titlefont $\bullet$ 20:38: transcribing previous notes on this}\smallskip
{\smallskip\noindent\bf\font\titlefont=cmr12 at 13pt\titlefont $\bullet$ 21:19: attemps to transcribe the figure}\smallskip
{\medskip\noindent\bf\font\titlefont=cmitt12 at 17pt\titlefont 2026-01-05}\medskip
{\smallskip\noindent\bf\font\titlefont=cmr12 at 13pt\titlefont $\bullet$ 21:25: block stack planning}\smallskip
{\medskip\noindent\bf\font\titlefont=cmitt12 at 17pt\titlefont 2026-01-06}\medskip
{\smallskip\noindent\bf\font\titlefont=cmr12 at 13pt\titlefont $\bullet$ 09:51: build out 1:1 representation of my blockstack notes}\smallskip
A good first step towards digitizing this data is producing a digital equivalent of the ink.
\par

Let's think about data first. figures are not yet in dagzet. How do we get them in there. Plopping things in 'ink/figs' and then connecting is tempting, but it breaks my ink generator scripts because those just traverse the ink namespace.
\par

Maybe adding a "fig" attribute to ink? I can select the node, then append a fig attribute that points to the png file.
\par

From there, I can make another query that queries by page and produces the JSON needed for that page. Another process can generate the PDF associated with that.
\par

Here's a snag. I have a "pg" attribute for the page, which is separate from the ink node which also has the page. To query all nodes associated with a page, I query for this page?
\par

No that won't work. Pages aren't necessarily distinct.
\par

Querying needs to happen from a namespace then, we're back to that.
\par

I wonder if the mknotes.sh query can be changed to produce a JSON object instead of just a JSON array of rows. I'd like to add another query related to figures associated with the namespace. It would contain the node containing the figure path, the address, and the page the figure is associated with. I think that should be enough data to extract info for a single page. I'll use the same single query, just different renderers for the data.
\par

well, now the query has been reworked to have a 'nodes' object for all the nodes. A 'figs' key would come next
\par

figs are still not in the knowledge graph at this point. in the near future, multiple graph figures will be associated with a particular namespace. at the moment, figures are just associated with their page. there needs to be something else that groups figures with a particular namespace, similar to what the node transcriptions are doing.
\par

a possible way forward is to integrate with the existing transcription pipeline used to generate the knowledge graph. add another column to graphs.txt containing a figs.txt file for that graph. that file will contain the page name for the graph, and the node name to give it. In the transcription generate script, it checks to see if that file exists, and batch generates those figure files. It would also need to produce dagzet output, which would be appeneded to the generated dagzet file.
\par

might work? gotta go now
\par

{\medskip\noindent\bf\font\titlefont=cmitt12 at 17pt\titlefont 2026-01-08}\medskip
{\smallskip\noindent\bf\font\titlefont=cmr12 at 13pt\titlefont $\bullet$ 10:01: getting some inertia up to implement this bit array operator}\smallskip
Came up with the C API. I quickly realized that since this has multiple variables, I'd need to use a stack to make it fully addressable. Instead, I'm opting make a version where the arguments comfortably fit into the 8 nibbles of the 32-bit word register. Addressing is a unsigned offset relative to the cursor (3 nibbles, 12 bits, enough for the blockstack). Values can be 16 bits, so that's 4 nibbles. 1 nibbles is reserved for size (up to 16 bits, 1-16 no empty sizes allowed)
\par

{\smallskip\noindent\bf\font\titlefont=cmr12 at 13pt\titlefont $\bullet$ 11:30: implement initial bit array test words (xr and xw)}\smallskip
no functionality, I just want to make sure the bit extraction works. very informal printf stuff.
\par

getting roped into a refactor with iscmd
\par

words implemented, and they are calling the bit array C functions. Next up is to implement those functions.
\par

{\medskip\noindent\bf\font\titlefont=cmitt12 at 17pt\titlefont 2026-01-11}\medskip
{\smallskip\noindent\bf\font\titlefont=cmr12 at 13pt\titlefont $\bullet$ 16:19: implementing C array functions}\smallskip
I think my math was off calculating the bits.
\par

Tests pass on the first try. Luck? Did I do it right?
\par

{\smallskip\noindent\bf\font\titlefont=cmr12 at 13pt\titlefont $\bullet$ 17:44: doing some interactive work to try some things out}\smallskip
I don't think I'm able to read/write a single bit, or anything that falls within one word.
\par

{\medskip\noindent\bf\font\titlefont=cmitt12 at 17pt\titlefont 2026-01-14}\medskip
{\smallskip\noindent\bf\font\titlefont=cmr12 at 13pt\titlefont $\bullet$ 10:24: refactoring the blockstack}\smallskip
refreshing myself on the current blockstack implementation. I'm pretty sure it's a linked list in an array sort of deal already. There must be a specific slot with the head of the list as well.
\par

Yup. It's a linked-list-in-an-array sort of deal. Let's call this array $L$, with element $n$ of the array being referred to as $L_n$. $L_0$ is the first item of the array, and it contains the address of the head of the list, $h$ initialized $h = 0$ (null). Then, the list gets initialized in the following way: $L_i \leftarrow h, h \leftarrow i$ for all $1 \le i \le 2^{10}$. So for $L_i = x$, $i$ refers to the address of the node if it's greater than 0 (with 0 being the head), and the value points to the next value in the linked list.
\par

I'm going to keep the addressing scheme the same with 0 being NULL. This leaves 1023 blocks max. Probably an okay limitation? Down the line, I'll want to have an externally managed chunk of memory that's 1024 blocks. I lose a block with this limitation.
\par

Writing out local "array 10" routines, for 10 bit arrays using the bit array routines I made.
\par

setting up the init/push/pop routines. seeing what fails, if there's failure.
\par

the tests are now failing in ways I understand. only the blocklist test is affected, and it has to do with the off-by-one addressing. I could fix the issue by changing the test values and be done with it, but I'm left wondering: should I change the address ordering to be ascending instead of descending? That would mean that no matter the size, the first element is 1 instead of 1023 or 1024.
\par

I was getting a weird return value of "3FF" from the original test. this is happening because I'm doing error checking, and the register isn't being overwritten. I do not have an error flag set up in place.
\par

Weird bug related to freeing blocks, a push operation. I'm doing a check to make sure the block isn't already free, which is to check to see if the node points to null (0) or not. A value that doesn't point to null is assumed to be in the list. When I made things initialized in reverse order (prev?), the last item (1013) points to 0, which I believe is correct. So, I might mark things popped off the list  as $L_i \leftarrow i$, which I believe works.
\par

tests refactored to work as expected. First item popped is 1. this all makes sense to me.
\par

{\smallskip\noindent\bf\font\titlefont=cmr12 at 13pt\titlefont $\bullet$ 12:13: I think I'm done? double checking my notes}\smallskip
{\smallskip\noindent\bf\font\titlefont=cmr12 at 13pt\titlefont $\bullet$ 12:15: Not done yet. address mapping is still needed}\smallskip
A block returned from this blocklist needs to be converted to an address in memory somehow. An address like {\tt 100.0F} should also work to, used to indicate at block {\tt 100} (256) plus word offset {\tt 0F} (15).
\par

Will think on the best approach for implementing.
\par

{\medskip\noindent\bf\font\titlefont=cmitt12 at 17pt\titlefont 2026-01-15}\medskip
{\smallskip\noindent\bf\font\titlefont=cmr12 at 13pt\titlefont $\bullet$ 08:41: address conversion}\smallskip
silly concerns: are 'bm' and 'mb' taken? That's block index to memory and memory to block.
\par

nice! okay those are the words.
\par

I messed up the test case. I need to multiply to get the memory address and THEN apply the bias.
\par

it all works now. tedious to get the tests right. might have to do with my ed workflow. it's still not quite as smooth as my established vim workflow
\par

{\smallskip\noindent\bf\font\titlefont=cmr12 at 13pt\titlefont $\bullet$ 10:31: I think we are done}\smallskip
{\medskip\noindent\bf\font\titlefont=cmitt12 at 17pt\titlefont 2026-01-16}\medskip
{\smallskip\noindent\bf\font\titlefont=cmr12 at 13pt\titlefont $\bullet$ 10:27: I think my offset calculations are wrong, nitty gritty details of memory addressing}\smallskip
At the very least, I believe there is an off by 1 error because block ids returned start at 1. Say the blocklist starts at zero. That's 5 blocks, or 320 words. block 1 should be at position 320, or 140. What is being returned right now is 0x180. a 0x40 difference means a block is being skipped due to the off by one. solution is subtract 1 from the block id before computing.
\par

This change makes it explicit so that addresses returned by the blockstack are 1-indexed and can never return 0.
\par

alternatively, maybe the blockstack can return zero-indexed values that are referenced by adding 1 to them. I'll have to think on this.
\par

second issue I'm going to double check are block bounds, something I think I was rambling about yesterday. If the blocklist starts at 0x00, the address of the first block should be 0x140. if the address is instead at 0x01, what should the address be at? Relatively speaking, the list take 5 blocks, but block addresses are always memory aligned. When the blocklist starts at 0x00, it begins at block 0 and ends at block 4. When the blocklist starts at 0x01, most of it is contained between blocks 0 and 4, but the last word spills out into block 5, thus containing 6 blocks. So, if blocks are memory-aligned, the address of the first usable block would skip block 5 (0x140) and go onto block 6 (0x180).
\par

At the time of writing, {\tt \#0 go \#1 bmpr } and {\tt \#1 go \#1 bmpr } both print the same value, which seems incorrect to me.
\par

Surprised with how precise I have to be about this and how easy it is to mess things up.
\par

\xrdef{node17}
{\medskip\noindent\bf\font\titlefont=cmss17 at 20pt\titlefont ink}\medskip
{\medskip\noindent\bf\font\titlefont=cmitt12 at 17pt\titlefont 2025-12-21}\medskip
{\smallskip\noindent\bf\font\titlefont=cmr12 at 13pt\titlefont $\bullet$ 09:34: sonilo now has an ink node}\smallskip
{\smallskip\noindent\bf\font\titlefont=cmr12 at 13pt\titlefont $\bullet$ 09:44: re-binning pages to be in sonilo ink}\smallskip
\xrdef{node18}
{\medskip\noindent\bf\font\titlefont=cmss17 at 20pt\titlefont initial parser}\medskip
{\medskip\noindent\bf\font\titlefont=cmitt12 at 17pt\titlefont 2025-12-28}\medskip
{\smallskip\noindent\bf\font\titlefont=cmr12 at 13pt\titlefont $\bullet$ 11:14: begin work implementing this graph drawing algo}\smallskip
Lots of time with pencil and paper these past few days, working out a way to transcribe these knowledge graphs I've drawing by hand.
\par

I think I managed to figure out a way to get a low-level representation of a graph using a system of weighted edges connecting quad trees.
\par

A drawn graph is made up of vertices that are quad trees. Each vertice has 4 ports labeled north, south, east, west. Visually, these correspond to the cardinal directions.
\par

An edge is used to connect two vertices together. There are two types of edges: horiztonal and vertical.
\par

A horizontal edge creates a directed edge from one nodes east port to another nodes west port. A vertical edge goes south to north. This makes graphs flow left to right, top down.
\par

An edge has a weight, which indicates how much space to add between the node. The amount to move is some kind of unit, like a tile.
\par

Defined edges are primarily used to draw connections between two vertices, but they can also be used for alignment and traversal purposes. In these cases, the edge can be marked to be invisible.
\par

Vertices can be of two types: a node or a junction. A node is the thing that contains the actual bit of information in the knowledge graph. The junction is used to make right angle bends for cables.
\par

{\smallskip\noindent\bf\font\titlefont=cmr12 at 13pt\titlefont $\bullet$ 11:38: some initial parser work}\smallskip
{\smallskip\noindent\bf\font\titlefont=cmr12 at 13pt\titlefont $\bullet$ 14:05: Okay, I have made an adjacency list, next steps}\smallskip
It's pretty cool to see what I wrote down on paper show up properly on the computer.
\par

What needs to happen next is traversal and drawing. Traversal might be a little bit faster to get to first, I can just print drawing commands to standard output.
\par

Needed: a queue, a "visited" list. Queue needs to take in a cursor position and pointer to a vertice.
\par

{\smallskip\noindent\bf\font\titlefont=cmr12 at 13pt\titlefont $\bullet$ 14:13: make a "visited" list}\smallskip
I think I'm just going to add an entry to the vertice list.
\par

{\smallskip\noindent\bf\font\titlefont=cmr12 at 13pt\titlefont $\bullet$ 14:16: onto a queue}\smallskip
{\smallskip\noindent\bf\font\titlefont=cmr12 at 13pt\titlefont $\bullet$ 14:40: back to queue}\smallskip
{\smallskip\noindent\bf\font\titlefont=cmr12 at 13pt\titlefont $\bullet$ 15:56: traversal works, the cursor offsets are wrong}\smallskip
Will need to troubleshoot later
\par

{\smallskip\noindent\bf\font\titlefont=cmr12 at 13pt\titlefont $\bullet$ 20:35: quick troubleshooting}\smallskip
I got it working! Drawing logic needs a bit more work. I am using weight to calculate the length of the line, but it doesn't account for the last little bits of the line. For junctions, the lines extend into the center point. For nodes, it might extend? Certainly there will be dots drawn to represent lines. So, maybe that's another drawing command, something that overlays stuff in node slots.
\par

Another thing that's kind of neat is that these draw commands can be run in arbitrary order. So, overlay commands can happen last.
\par

{\medskip\noindent\bf\font\titlefont=cmitt12 at 17pt\titlefont 2025-12-29}\medskip
{\smallskip\noindent\bf\font\titlefont=cmr12 at 13pt\titlefont $\bullet$ 08:48: adding overlay drawing commands}\smallskip
How I am imagining it will go: two overlays happen for every line. A line always connects two vertices together. Vertices can either be a junction (+) or a node (\$).
\par

The vertice type determines the kind of overlay to draw. Nodes will have dots drawn on the eastern and southern ports to indicate arrowheads (these are input ports), and they will possibly have some lines to draw as well, depending on how much whitespace there is. Junctions will have lines from the port to the center of the tile.
\par

Overlay commands will need to specify which direction to draw (N, S, E, W). As the name implies, the overlay commands should be able to write on top of previously drawn items in that tile.
\par

So, overlay data needed: position (tile X, Y), type (junction, node), and port to draw (N, S, E, W).
\par

This could be good pseudocode: OVERLAY: J 1 2 E
\par

{\smallskip\noindent\bf\font\titlefont=cmr12 at 13pt\titlefont $\bullet$ 09:33: overlays seem to look right. now to plan for drawing}\smallskip
I'm actually thinking about making another small program for drawing 1-bit graphics. I can then take the output commands from dagdraw and turn them into these lower level commands.
\par

The intermediate program could be an awk script
\par

So what will these commands consist of?
\par

Rectangles (outline and filled), hlines, vlines. I think that's all that's required for now (ignoring text).
\par

rectangle (rc)/filled rectangle (rf): x, y, w, h
\par

vline (lv): x, y, length. positive length is downward, negative is upward (positive only for now)
\par

hline (lh): x, y, length: positive length is eastward, negative westword (positive only for now)
\par

wr: write file to disk
\par

sz: size (w, h) of output
\par

{\smallskip\noindent\bf\font\titlefont=cmr12 at 13pt\titlefont $\bullet$ 09:42: okay testing this out}\smallskip
{\smallskip\noindent\bf\font\titlefont=cmr12 at 13pt\titlefont $\bullet$ 11:06: "bitr" command language seems to work, writing initial command generator}\smallskip
{\smallskip\noindent\bf\font\titlefont=cmr12 at 13pt\titlefont $\bullet$ 11:31: VLINE/HLINE/NODE commands seem to work, onto OVERLAY}\smallskip
Uh oh, I tihnk my overlay code might be wrong. I don't think it's correctly calculating coordinates.
\par

Ah, I think it was an off-by-one thing with weights.
\par

I think the issues are more rotten than I expected. the junctions are sometimes off-by-one, sometimes not in the dagdraw output.
\par

Huh. Maybe it's just off-by-one? It's possible I might not be computing the overlays correctly. -- Rendering just west edges: J1w, J5w, J6w looks good. But hang on, there are other edges here that shoulnd't be here.
\par

Ah, typo when getting the type for "b". Let's see if that clears things up.
\par

Okay, that works now! phew.
\par

Onto node overlays now.
\par

It works! Amazing.
\par

{\smallskip\noindent\bf\font\titlefont=cmr12 at 13pt\titlefont $\bullet$ 12:37: going to add CLI args, test ex1}\smallskip
seems to work! I'll need to do more tweaks for things like configuring the initial origin, and calculating dimensions for the final output. But, it does what I want.
\par

\xrdef{node19}
{\medskip\noindent\bf\font\titlefont=cmss17 at 20pt\titlefont autoscaling}\medskip
{\medskip\noindent\bf\font\titlefont=cmitt12 at 17pt\titlefont 2025-12-29}\medskip
{\smallskip\noindent\bf\font\titlefont=cmr12 at 13pt\titlefont $\bullet$ 18:11: work on auto-scaling canvas size}\smallskip
I think the way I'll do this is to keep track of the max/min x and y coordinates, and then use that to determine the bounds.
\par

So, it'll be a two-pass traversal. Traverse it once to get the bounds, starting at (0, 0). Then run it again and add some kind of offset.
\par

well, it doesn't work yet.
\par

{\smallskip\noindent\bf\font\titlefont=cmr12 at 13pt\titlefont $\bullet$ 21:36: taking another swing at it}\smallskip
Hey, it works. The canvas will auto-adjust to fit the boundary.
\par

\xrdef{node20}
{\medskip\noindent\bf\font\titlefont=cmss17 at 20pt\titlefont fig prototype}\medskip
{\medskip\noindent\bf\font\titlefont=cmitt12 at 17pt\titlefont 2025-12-29}\medskip
{\smallskip\noindent\bf\font\titlefont=cmr12 at 13pt\titlefont $\bullet$ 22:00: figure out one level higher}\smallskip
Something that can generate the edge code a little bit easier
\par

I got my "arch" design, parallel, and sequence configurations working in a macro language I built in awk. More intuitive than I expectred.
\par

\xrdef{node21}
{\medskip\noindent\bf\font\titlefont=cmss17 at 20pt\titlefont labels}\medskip
{\medskip\noindent\bf\font\titlefont=cmitt12 at 17pt\titlefont 2025-12-30}\medskip
{\smallskip\noindent\bf\font\titlefont=cmr12 at 13pt\titlefont $\bullet$ 07:24: label support}\smallskip
The last thing I need to get this actually usuable is to add labels for nodes.
\par

I don't really want strings in bitr, so I think I'm going to stick with the 6-bit encoding system I used in dagdraw to encode the text.
\par

The bitmap font will be embedded in the binary. I've selected a few to try out. 8x8 fonts.
\par

rectangles are going to need to be stretched a bit to fit the text. I'm thinking I'll make it a 2 character label max for now (that's ~1300 nodes in base 36, which I use for indexed nodes. plenty for this scale).
\par

{\smallskip\noindent\bf\font\titlefont=cmr12 at 13pt\titlefont $\bullet$ 07:32: stretch the grid size in draw.awk}\smallskip
Make it so 2 characters can fit.
\par

Okay that mostly seems to work. Took some fiddly recalculations. pixel accurate centering takes a bit of thought. Lots of odd integers.
\par

{\smallskip\noindent\bf\font\titlefont=cmr12 at 13pt\titlefont $\bullet$ 08:30: get a pixel font working... with centering?}\smallskip
Pristine seems to be the current winner.
\par

I got a glyph printed, but I need to adjust spacing for the boxes because the letters won't fit.
\par

{\smallskip\noindent\bf\font\titlefont=cmr12 at 13pt\titlefont $\bullet$ 08:37: initial text rendering works!}\smallskip
This is good enough for display purposes. However, the text ideally should be horizontally centered.
\par

And we are centered. Excellent.
\par

I'll need to work out the "fig" language a bit more, but the fundamentals are in place to transcribe these knowledge graphs.
\par

{\smallskip\noindent\bf\font\titlefont=cmr12 at 13pt\titlefont $\bullet$ 09:13: making the boxes a little bit bigger}\smallskip
\xrdef{node22}
{\medskip\noindent\bf\font\titlefont=cmss17 at 20pt\titlefont duplicate node}\medskip
{\medskip\noindent\bf\font\titlefont=cmitt12 at 17pt\titlefont 2025-12-31}\medskip
{\smallskip\noindent\bf\font\titlefont=cmr12 at 13pt\titlefont $\bullet$ 12:22: Duplicate node support}\smallskip
Needed for a graph I was trying to transcribe yesterday.
\par

{\smallskip\noindent\bf\font\titlefont=cmr12 at 13pt\titlefont $\bullet$ 13:11: Working out a strategy}\smallskip
The strategy I have in mind is to denote duplicate nodes in the following way: {\tt \$XY.A} That is to say, there is a node XY and this has a unique identifier of "A". When the drawing commands get produced, this will go away.
\par

\xrdef{node23}
{\medskip\noindent\bf\font\titlefont=cmss17 at 20pt\titlefont butlp}\medskip
{\medskip\noindent\bf\font\titlefont=cmitt12 at 17pt\titlefont 2026-01-02}\medskip
{\smallskip\noindent\bf\font\titlefont=cmr12 at 13pt\titlefont $\bullet$ 15:48: this is a doable self-contained task}\smallskip
It'll mostly be copy-paste work from sndkit.
\par

\xrdef{node24}
{\medskip\noindent\bf\font\titlefont=cmss17 at 20pt\titlefont notes integration}\medskip
{\medskip\noindent\bf\font\titlefont=cmitt12 at 17pt\titlefont 2026-01-02}\medskip
{\smallskip\noindent\bf\font\titlefont=cmr12 at 13pt\titlefont $\bullet$ 16:16: dagdraw notes integration: initial thoughts}\smallskip
I'm building a TeX generator for viewing my transcribed reading notes, and I think it would make sense to include these generated figures.
\par

To best read the figure, it needs an index. Ideally, the figure and index should be on the same page. These typically take up a page written by hand, so this shouldn't be too difficult.
\par

figures need to have the same pixel ratio, otherwise things will look too stretched out. Generate some kind of metrics file? That way, I won't need to parse the image files in the tex generator.
\par

the index doesn't necessarily need to have all the text, just the ID, the abbreviation, and page number.
\par

How to get figures into the note? They need to get imported into the dag. Maybe a figs namespace? They are still tied to ink IDs, so maybe the figs are located in ink, and then connected to a figs source node and the ink.
\par

A notes generator would have some way of traversing the child nodes of some figs namespace.
\par

{\medskip\noindent\bf\font\titlefont=cmitt12 at 17pt\titlefont 2026-01-06}\medskip
{\smallskip\noindent\bf\font\titlefont=cmr12 at 13pt\titlefont $\bullet$ 09:51: build out 1:1 representation of my blockstack notes}\smallskip
A good first step towards digitizing this data is producing a digital equivalent of the ink.
\par

Let's think about data first. figures are not yet in dagzet. How do we get them in there. Plopping things in 'ink/figs' and then connecting is tempting, but it breaks my ink generator scripts because those just traverse the ink namespace.
\par

Maybe adding a "fig" attribute to ink? I can select the node, then append a fig attribute that points to the png file.
\par

From there, I can make another query that queries by page and produces the JSON needed for that page. Another process can generate the PDF associated with that.
\par

Here's a snag. I have a "pg" attribute for the page, which is separate from the ink node which also has the page. To query all nodes associated with a page, I query for this page?
\par

No that won't work. Pages aren't necessarily distinct.
\par

Querying needs to happen from a namespace then, we're back to that.
\par

I wonder if the mknotes.sh query can be changed to produce a JSON object instead of just a JSON array of rows. I'd like to add another query related to figures associated with the namespace. It would contain the node containing the figure path, the address, and the page the figure is associated with. I think that should be enough data to extract info for a single page. I'll use the same single query, just different renderers for the data.
\par

well, now the query has been reworked to have a 'nodes' object for all the nodes. A 'figs' key would come next
\par

figs are still not in the knowledge graph at this point. in the near future, multiple graph figures will be associated with a particular namespace. at the moment, figures are just associated with their page. there needs to be something else that groups figures with a particular namespace, similar to what the node transcriptions are doing.
\par

a possible way forward is to integrate with the existing transcription pipeline used to generate the knowledge graph. add another column to graphs.txt containing a figs.txt file for that graph. that file will contain the page name for the graph, and the node name to give it. In the transcription generate script, it checks to see if that file exists, and batch generates those figure files. It would also need to produce dagzet output, which would be appeneded to the generated dagzet file.
\par

might work? gotta go now
\par

{\medskip\noindent\bf\font\titlefont=cmitt12 at 17pt\titlefont 2026-01-07}\medskip
{\smallskip\noindent\bf\font\titlefont=cmr12 at 13pt\titlefont $\bullet$ 09:26: continuing where I left off yesterday}\smallskip
The "fig" attribute is successfully showing up in the database. I still am going to need dimensions.
\par

opting for imagemagick to get this. I'm adding them as attributes to the fig node.
\par

I had to massage the SQLite query a bit to produce the JSON output I wanted, but the info is in there for both figures and nodes.
\par

now to make a test generator
\par

working inkpage. mostly satisfied with this
\par

Note that this is hardcoded to be only one page. What I'll need to do later is follow-up and make it so it generates a PDF for all pages
\par

{\medskip\noindent\bf\font\titlefont=cmitt12 at 17pt\titlefont 2026-01-11}\medskip
{\smallskip\noindent\bf\font\titlefont=cmr12 at 13pt\titlefont $\bullet$ 14:30: let's see if my new eps backend works okay on linux}\smallskip
seems to work.
\par

going to write down a few notes for myself for future me. there are a lot of small steps.
\par

\xrdef{node25}
{\medskip\noindent\bf\font\titlefont=cmss17 at 20pt\titlefont eps}\medskip
{\medskip\noindent\bf\font\titlefont=cmitt12 at 17pt\titlefont 2026-01-10}\medskip
{\smallskip\noindent\bf\font\titlefont=cmr12 at 13pt\titlefont $\bullet$ 14:47: fiddling with eps output using cairo}\smallskip
I got it looking fairly decent, but there's a small issue with overlapping lines. I think it has to do with how my system wants to draw line segements individually.
\par

I think to improve this, I'd want to add a step that combines contiguous line segments together.
\par

Anyways, time to see if my ink renderer gets improved with eps output.
\par

\xrdef{node26}
{\medskip\noindent\bf\font\titlefont=cmss17 at 20pt\titlefont path drawing}\medskip
{\medskip\noindent\bf\font\titlefont=cmitt12 at 17pt\titlefont 2026-01-10}\medskip
{\smallskip\noindent\bf\font\titlefont=cmr12 at 13pt\titlefont $\bullet$ 19:45: DAGDraw path drawing algorithm}\smallskip
Musing possible approaches to improving the path drawing logic. in the vectorized output of DAGDraw. The main process involves grouping segments together into contiguous lines to minimize the number of weird overlapping lines that can show up in output.
\par

\xrdef{node27}
{\medskip\noindent\bf\font\titlefont=cmss17 at 20pt\titlefont gesture}\medskip
{\medskip\noindent\bf\font\titlefont=cmitt12 at 17pt\titlefont 2026-01-11}\medskip
{\smallskip\noindent\bf\font\titlefont=cmr12 at 13pt\titlefont $\bullet$ 20:23: everything is a gesture}\smallskip
Thinking ahead about high-level concepts. This music environment I want to make is going to involve the concept of Gesture, so what would happen if I go to an extreme and make everything a Gesture. Gesture as a core primitive. Composing with gesture.
\par

\xrdef{node28}
{\medskip\noindent\bf\font\titlefont=cmss17 at 20pt\titlefont fix block navigation}\medskip
{\medskip\noindent\bf\font\titlefont=cmitt12 at 17pt\titlefont 2026-01-15}\medskip
{\smallskip\noindent\bf\font\titlefont=cmr12 at 13pt\titlefont $\bullet$ 15:32: 'bw' and 'gb' need to be fixed}\smallskip
bw: is helpful function that needs to be either removed or reworked. It needs to factor in the cursor position, which will contain the address of the block stack.
\par

gb: goto block. this should eventually be reworked to use the 'bm' word like so: {\tt \#01 bm go}.
\par

I'm only using it in tests for now, so nothing is going to fall apart. I just can't use them anymore and I should be fine. Deprecated.
\par

{\medskip\noindent\bf\font\titlefont=cmitt12 at 17pt\titlefont 2026-01-16}\medskip
{\smallskip\noindent\bf\font\titlefont=cmr12 at 13pt\titlefont $\bullet$ 10:27: I think my offset calculations are wrong, nitty gritty details of memory addressing}\smallskip
At the very least, I believe there is an off by 1 error because block ids returned start at 1. Say the blocklist starts at zero. That's 5 blocks, or 320 words. block 1 should be at position 320, or 140. What is being returned right now is 0x180. a 0x40 difference means a block is being skipped due to the off by one. solution is subtract 1 from the block id before computing.
\par

This change makes it explicit so that addresses returned by the blockstack are 1-indexed and can never return 0.
\par

alternatively, maybe the blockstack can return zero-indexed values that are referenced by adding 1 to them. I'll have to think on this.
\par

second issue I'm going to double check are block bounds, something I think I was rambling about yesterday. If the blocklist starts at 0x00, the address of the first block should be 0x140. if the address is instead at 0x01, what should the address be at? Relatively speaking, the list take 5 blocks, but block addresses are always memory aligned. When the blocklist starts at 0x00, it begins at block 0 and ends at block 4. When the blocklist starts at 0x01, most of it is contained between blocks 0 and 4, but the last word spills out into block 5, thus containing 6 blocks. So, if blocks are memory-aligned, the address of the first usable block would skip block 5 (0x140) and go onto block 6 (0x180).
\par

At the time of writing, {\tt \#0 go \#1 bmpr } and {\tt \#1 go \#1 bmpr } both print the same value, which seems incorrect to me.
\par

Surprised with how precise I have to be about this and how easy it is to mess things up.
\par

{\medskip\noindent\bf\font\titlefont=cmitt12 at 17pt\titlefont 2026-02-07}\medskip
{\smallskip\noindent\bf\font\titlefont=cmr12 at 13pt\titlefont $\bullet$ 11:05: bw rewrite scoping broad strokes}\smallskip
For blockword to work again, it would need to have the cursor set to be the global blockstack. from there, it computes the word address of the first block, then adds the offset.
\par

I also need to update the tests using this (which I don't think are many) to ensure that the global blocks are pointed at the cursor.
\par

{\smallskip\noindent\bf\font\titlefont=cmr12 at 13pt\titlefont $\bullet$ 11:08: 'bm' is off by one, more or less}\smallskip
I don't have this in my head anymore.
\par

\xrdef{node29}
{\medskip\noindent\bf\font\titlefont=cmss17 at 20pt\titlefont news}\medskip
{\medskip\noindent\bf\font\titlefont=cmitt12 at 17pt\titlefont 2026-01-22}\medskip
{\smallskip\noindent\bf\font\titlefont=cmr12 at 13pt\titlefont $\bullet$ 10:13: stretching sonilo task days}\smallskip
We're 67 days behind now.
\par

Well, now it says end of july things might be done. Not gonna happen probably. I just kinda need this memory stuff to be behind me.
\par

{\medskip\noindent\bf\font\titlefont=cmitt12 at 17pt\titlefont 2026-02-03}\medskip
{\smallskip\noindent\bf\font\titlefont=cmr12 at 13pt\titlefont $\bullet$ 09:22: delaying sonilo about 50 days to keep the backlog down}\smallskip
My projected deadline says july 3rd, 2026, now. Probably not realistic, considering how I'm only working one day a week on this.
\par

{\medskip\noindent\bf\font\titlefont=cmitt12 at 17pt\titlefont 2026-02-07}\medskip
{\smallskip\noindent\bf\font\titlefont=cmr12 at 13pt\titlefont $\bullet$ 11:02: review. planning my day}\smallskip
Last week: block execution, scoping out ugens and subtasks.
\par

This week: get reference counting structure implemented.
\par

{\medskip\noindent\bf\font\titlefont=cmitt12 at 17pt\titlefont 2026-02-14}\medskip
{\smallskip\noindent\bf\font\titlefont=cmr12 at 13pt\titlefont $\bullet$ 09:53: planning the day}\smallskip
Looks like I finished the reference counting structure last week.
\par

Agenda has bulk instructions, block navigation fixes, and the ugen stack.
\par

Block navigation fixes can wait.
\par

{\medskip\noindent\bf\font\titlefont=cmitt12 at 17pt\titlefont 2026-02-15}\medskip
{\smallskip\noindent\bf\font\titlefont=cmr12 at 13pt\titlefont $\bullet$ 10:28: adding block pool task}\smallskip
\xrdef{node30}
{\medskip\noindent\bf\font\titlefont=cmss17 at 20pt\titlefont instruction primitive}\medskip
{\medskip\noindent\bf\font\titlefont=cmitt12 at 17pt\titlefont 2026-01-24}\medskip
{\smallskip\noindent\bf\font\titlefont=cmr12 at 13pt\titlefont $\bullet$ 14:10: initial words}\smallskip
See "ugen bytecode" task for previous thoughts on this.
\par

goal: execute an instruction given a word.
\par

The instruction format divides a word into two parts: data (16 bits) and command (16 bits).
\par

For nodes, the data portion is most likely going to be a memory address, though for this task it can just be treated as a literal constant.
\par

The command portion is going to need some thought. 16-bits is more instructions than we'll ever use (I think?). Using it as a lookup table could be nice.
\par

16 bits can encode some kind of alphabetic or alphanumeric string. I know I can atleast get three alphabetic or alphanumeric characters. If I know that numbers are usually suffixes, I might be able to pack 3 letters and 3 numbers like ABC123 into 16 bits.
\par

16 bits pack 3 letters. 14 bits packs 3 letters (with constraints on letter choice) plus 2 bits for a number. Going to just stick with the 16 bit encoding approach.
\par

lookup table can be a linear probing dictionary.
\par

How to create an instruction? There needs to be a way to specify the data and also the command. Data will sometimes be a memory address, probably created first before calling the command. Entering in a command should effectively push bits onto the word.
\par

{\smallskip\noindent\bf\font\titlefont=cmr12 at 13pt\titlefont $\bullet$ 19:56: some initial scaffolding}\smallskip
Setting up structs and boilerplate. This glue code will exist outside the VM, since the routines will be implemented in C.
\par

Threw together some untested code for core dictionary operations. I've opted for linear probing. It's simple. No linked lists. The hash is modulo for now. It's fine. I can tweak it later.
\par

Encoding comes next. Strings need to be keys. I spent more time than I care to admit working out math looking into how to milk as many of the 16 bits as possible. With a mixed radix system, you can get 3 letters plus a value that can be 4 states (0, 1, 2, 3). Or, you can do 3 5-bit characters plus 1 bit (probably lose that bit). 5 bits gives you letters plus 6 specials (TBD? no plans yet), and is way simpler to build, so I'm going to go with that.
\par

Now comes the interesting bit of hooking this functionality into the test system. I'm going to use 'ex' for instruction execution. I'll have a test function called 'INC' that increments the data to be plus 1.
\par

Some new adjustments to the memwrite state machine required. Single quote (') will turn on 5-bit encoding. another single quote (') will turn it off.
\par

Interesting. I have to keep the bits nibble-aligned. So, I'm going to make the enclosing "'" shift the register 1 bit. assuming 3 characters, that should make it fit into 16 bits.
\par

And it works. I'm applying the 1-bit shift for the string to key utility. So, generally speaking, from this point on, I do declare that opcode names are exactly 3 letters (well, 3 5-bit 'symbols' which are an alphabet and 6 other symbols TBD).
\par

\xrdef{node31}
{\medskip\noindent\bf\font\titlefont=cmss17 at 20pt\titlefont bulk instructions}\medskip
{\medskip\noindent\bf\font\titlefont=cmitt12 at 17pt\titlefont 2026-01-24}\medskip
{\smallskip\noindent\bf\font\titlefont=cmr12 at 13pt\titlefont $\bullet$ 14:33: initial words}\smallskip
Implement the ability to render a block of instructions.
\par

Conditional jumping isn't going to be used here because DSP ugens just execute things in order.
\par

Probably utilize the existing array structure implemented, and figure out a decent way to append instructions to the array.
\par

{\medskip\noindent\bf\font\titlefont=cmitt12 at 17pt\titlefont 2026-01-31}\medskip
{\smallskip\noindent\bf\font\titlefont=cmr12 at 13pt\titlefont $\bullet$ 07:31: hoping to work on this today}\smallskip
{\smallskip\noindent\bf\font\titlefont=cmr12 at 13pt\titlefont $\bullet$ 16:12: bulk instructions: implementation}\smallskip
Setting up the initial unit test. The actual implementation is trivial.
\par

So this shouldn't be too strenuous. I'll use the exist block array, and append instructions to it.
\par

Cute idea: make the program be based on the konami code: up, up, down, down, left, right, left, right, B, A. up/down can be increase, decrease by one. left/right refers to bitshifting. B is bitwise negation. A is reverse bits. I'll add some print opcodes as well.
\par

hoping that 'bx' isn't taken. that would be a good acronym to use.
\par

once again, my left/right confusion has worked against me. I mixed up left shift and right shift operations which made my expected results wrong. Also, I forgot what the binary nibble of 6 was. Off-by-one.
\par

'bx' has been implemented.
\par

side note: these instruction operations should eventually be consolidated into C functions, not just inline code.
\par

{\medskip\noindent\bf\font\titlefont=cmitt12 at 17pt\titlefont 2026-02-14}\medskip
{\smallskip\noindent\bf\font\titlefont=cmr12 at 13pt\titlefont $\bullet$ 09:59: This is basically complete. Just wrap things in functions}\smallskip
I left some TODOs for myself. Trivial work.
\par

{\smallskip\noindent\bf\font\titlefont=cmr12 at 13pt\titlefont $\bullet$ 20:37: addressing the small TODOS}\smallskip
\xrdef{node32}
{\medskip\noindent\bf\font\titlefont=cmss17 at 20pt\titlefont ugens}\medskip
{\medskip\noindent\bf\font\titlefont=cmitt12 at 17pt\titlefont 2026-01-24}\medskip
{\smallskip\noindent\bf\font\titlefont=cmr12 at 13pt\titlefont $\bullet$ 14:44: getting close to needing to think about this}\smallskip
I am currently scoping out the ugen bytecode task, which has been greatly simplified over the months. What was once a whole VM has now turned into a sequential list of operations that read/write to data stored at a particular memory address.
\par

This simpler instruction-based format means some more responsibilities for the ugen API. Mainly, I/O stuff. Nodes need inputs and nodes produce outputs.
\par

Luckily, I have some pre-existing work for this: patchwerk. It's looking like I'm going to be building out something very similar to this.
\par

These inputs and outputs will have their addresses baked into the data somehow.
\par

An instruction takes in one data parameter. Perhaps this is a structure containing ugen data. Ugens have an array of input params (blocks, constants), an array of outputs (blocks, probably not constants), and data pointing to the internal state of the DSP.
\par

Audio rate signals are blocks. Patchwerk uses a fixed size buffer (block) pool and a buffer (block) stack with reference counting to manage blocks. A similar kind of reference counting system could possibly be used.
\par

When it comes to memory management, I'll take simple yet tedious/error-prone over complex yet frictionless. There are going to be abstractions on top of this, especially when gesture becomes more involved (leaning towards a kind of system where everything can be expressed as a gesture).
\par

{\medskip\noindent\bf\font\titlefont=cmitt12 at 17pt\titlefont 2026-01-26}\medskip
{\smallskip\noindent\bf\font\titlefont=cmr12 at 13pt\titlefont $\bullet$ 19:30: Ugen Structure}\smallskip
Elaborated on ideas related to the structure of a ugen, and some of the operations used to create them. A patchwerk inspired reference counting system is what I'm thinking about. The change I'm making is when garbage collection happens for blocks. The change will make it so the sets of input and output blocks will be disjoint, making DSP code a bit easier to work with.
\par

{\medskip\noindent\bf\font\titlefont=cmitt12 at 17pt\titlefont 2026-01-31}\medskip
{\smallskip\noindent\bf\font\titlefont=cmr12 at 13pt\titlefont $\bullet$ 07:29: ugen subtasks can be scoped out now I think}\smallskip
{\smallskip\noindent\bf\font\titlefont=cmr12 at 13pt\titlefont $\bullet$ 11:05: ugen structure 2}\smallskip
Taking another stab at outlining the componets that make up unit generators in an effort to see what needs to be built bottoms up
\par

{\smallskip\noindent\bf\font\titlefont=cmr12 at 13pt\titlefont $\bullet$ 13:18: Ugen prelim(inaries)}\smallskip
Having previously outlined top-down the components that make up ugens, this graph, examines at the leaves of that graph (the preliminary components needed to get ugens working), and pulls out more details and relationships between them.
\par

Ports, ref counters, and stacks, in approximately that order.
\par

ref counters, stacks, and then ports, in approximately that order.
\par

{\smallskip\noindent\bf\font\titlefont=cmr12 at 13pt\titlefont $\bullet$ 14:34: creating new tasks for ugens, writing words}\smallskip
{\smallskip\noindent\bf\font\titlefont=cmr12 at 13pt\titlefont $\bullet$ 14:55: Garbage collection after ports, order matters}\smallskip
I may have mentioned this already. But it's one of those things I can see slipping from my mind. So, we're writing it down again.
\par

Unit generators need to be able to share and reuse blocks. A stack and a reference counting system is used to help with this.
\par

A ugen has a certain number of input ports and output ports. A ugen generator will allocate data required for the DSP state, set up the ports, and then finally do garbage collection on the reference counter.
\par

A pop operation decreases reference count associated with a particular memory address of a block. A push operation will increase it. A garbage collect operation will mark that block as available for re-use.
\par

The garbage collection happens last, after all stack operations. Any time a new output parameter is pushed onto the stack, it will be unable to claim a recently popped block used by the input, guaranteeing that input/output blocks will be disjoing. Making the output/input blocks disjoint prevents things like accidental overwrites to happen.
\par

{\smallskip\noindent\bf\font\titlefont=cmr12 at 13pt\titlefont $\bullet$ 16:01: task ordering updates, ugens child of some other tasks now}\smallskip
ilo interpreter stuff needs to at least happen after ugen stuff. And it's a prereq for any gesture stuff.
\par

By the time unit generators are finished, I think thinking about using a higher-level abstraction like Konilo may be warranted. At that point, it'll start to feel possible to build demos.
\par

\xrdef{node33}
{\medskip\noindent\bf\font\titlefont=cmss17 at 20pt\titlefont ugen ports}\medskip
{\medskip\noindent\bf\font\titlefont=cmitt12 at 17pt\titlefont 2026-01-31}\medskip
{\smallskip\noindent\bf\font\titlefont=cmr12 at 13pt\titlefont $\bullet$ 14:34: creating new tasks for ugens, writing words}\smallskip
{\smallskip\noindent\bf\font\titlefont=cmr12 at 13pt\titlefont $\bullet$ 14:43: initial words}\smallskip
a unit generator does I/O via things called "ports". A port can either be a pointer to a block (an audio rate signal), or it can be a constant.
\par

The port implementation is designed to fit into a single word. The first 2 bits declare the type of port. The remaining 30 bits are used to store data.
\par

one type of port is a block. A block contains a 16-bit word address to the start of a block.
\par

Another type of port is a constant, which is encoded in 30 bits. The format I'm going with here is a signed fixed point 16.14 number format. I think that should be enough resolution for this context.
\par

Unit generators will call DSP routines, and so it's probably a good idea to turn a port into an equivalent C-struct. Here, the constants would be converted into C floats, and the word addresses would be turned into pointers.
\par

\xrdef{node34}
{\medskip\noindent\bf\font\titlefont=cmss17 at 20pt\titlefont ugen refcount}\medskip
{\medskip\noindent\bf\font\titlefont=cmitt12 at 17pt\titlefont 2026-01-31}\medskip
{\smallskip\noindent\bf\font\titlefont=cmr12 at 13pt\titlefont $\bullet$ 14:34: creating new tasks for ugens, writing words}\smallskip
{\smallskip\noindent\bf\font\titlefont=cmr12 at 13pt\titlefont $\bullet$ 14:55: Garbage collection after ports, order matters}\smallskip
I may have mentioned this already. But it's one of those things I can see slipping from my mind. So, we're writing it down again.
\par

Unit generators need to be able to share and reuse blocks. A stack and a reference counting system is used to help with this.
\par

A ugen has a certain number of input ports and output ports. A ugen generator will allocate data required for the DSP state, set up the ports, and then finally do garbage collection on the reference counter.
\par

A pop operation decreases reference count associated with a particular memory address of a block. A push operation will increase it. A garbage collect operation will mark that block as available for re-use.
\par

The garbage collection happens last, after all stack operations. Any time a new output parameter is pushed onto the stack, it will be unable to claim a recently popped block used by the input, guaranteeing that input/output blocks will be disjoing. Making the output/input blocks disjoint prevents things like accidental overwrites to happen.
\par

{\smallskip\noindent\bf\font\titlefont=cmr12 at 13pt\titlefont $\bullet$ 15:06: initial words}\smallskip
Build a reference counting system. Essentially, it'll be a list of block addresses and a count, indicating how many times they'll be used.
\par

Initially, there will need to be an operation that adds new addresses to the list. This will find an available slot and store it there.
\par

There needs to be a way to delete an address from the list, clearing it and making the slot available to be overriden by another add operation.
\par

A "find" operation will look up an address to see if it associated with the RC list. This will be helpful for other operations.
\par

Counts need to increase/decrease with push/pop stack operations.
\par

Sometimes blocks need to be held indefinitely. These will be called "hold" and "unhold" operations. A "hold" bit will need to be made.
\par

A "collect" operation is used to sweep through the list, finding anything with a used slot with a count of 0, and mark it as unused. A "used" bit will need to be made.
\par

Memory address + 2 bits for used/hold = 18 bits for a slot. Maybe pack this along with the stack?
\par

{\smallskip\noindent\bf\font\titlefont=cmr12 at 13pt\titlefont $\bullet$ 15:19: packing the stack and refcounter together using some arithmetic}\smallskip
One refcount entry needs 18 bits of data. 64 slots is 1152 bits of data, conveniently 36 words of data. 28 words left for the stack.
\par

slot indexes need 6 bits of data. If the stack uses an array, it's possible to have upper limit of approximiately 149 for a stack size. Round to the nearest power of 2, 128. 768 bits, 24 words. But, that doesn't include the stack pointer, which would need to be at least 8 bits to include 0.
\par

Alternatively, use an array structure, and have the first entry contain the stack pointer (the number of items on the stack). a 6-bit number would have a range 0-63, or a stack size of 63. a 7-bit number would have a range 0-127, or a stack size of 127.
\par

a 6-bit stack needs an array of 64 6-bit numbers, 12 words. A 7-bit stack needs an array of 128 7-bit numbers, 28 words.
\par

It turns out that 64 18-bit numbers plus 128 7-bit numbers yields 32 64-bit numbers, or one block. How lovely.
\par

{\medskip\noindent\bf\font\titlefont=cmitt12 at 17pt\titlefont 2026-02-07}\medskip
{\smallskip\noindent\bf\font\titlefont=cmr12 at 13pt\titlefont $\bullet$ 11:09: attempting to accomplish this today}\smallskip
Based on my logs, it looks like I've already scoped out the functionality of this.
\par

Looks like I haven't actually implemented anything yet.  Guess that's what I'll be looking into today.
\par

Reference count list: each entry 18 bits. Will take up first part of a block (the remaining bits will be reserved for the stack).
\par

The global bit array will come in handy here. I'll build things using that.
\par

operations: add, delete, find, hold, unhold, collect.
\par

{\smallskip\noindent\bf\font\titlefont=cmr12 at 13pt\titlefont $\bullet$ 12:10: refcount list brainstorm (part 1)}\smallskip
Structuring thoughts on how to design and implement this reference counting structure. I broke it down into data an operations, and got really sucked into the data aspect of it. I wanted to tune things so that I could optimimally fit a stack and a refcount list on a single block of memory.
\par

{\smallskip\noindent\bf\font\titlefont=cmr12 at 13pt\titlefont $\bullet$ 14:27: refcount list brainstorm (part 2)}\smallskip
Outlined the core operations, starting drilling into some of them. Also started thinking about some of the helper utilites.
\par

{\smallskip\noindent\bf\font\titlefont=cmr12 at 13pt\titlefont $\bullet$ 15:15: refcount list brainstorm (part 2)}\smallskip
more fleshing out details for each of the operations.
\par

{\smallskip\noindent\bf\font\titlefont=cmr12 at 13pt\titlefont $\bullet$ 17:06: refcount list (brainstorm) (part 2) (continued)}\smallskip
I returned to the data problem. With my overengineered equation on optimum size in hand, I started plugging in some variables to see how they'd turn out. I tried a 7-bit stack with 64 slots, then made it 48 slots.
\par

{\smallskip\noindent\bf\font\titlefont=cmr12 at 13pt\titlefont $\bullet$ 17:50: packing calculations were wrong}\smallskip
The size of an entry is not 18 bytes, that's only for the flags and memory address. A counter needs to be supplied as well. I think 4 bits ought to be enough (well before 15 does it make sense to use the hold mechanism instead).
\par

According to my calculations (and yes, I did calculations): 48 entries of 22 bits, a stack size of 96 with 7-bit ints. With all that, you get about 9 words left in the block.
\par

{\smallskip\noindent\bf\font\titlefont=cmr12 at 13pt\titlefont $\bullet$ 17:54: write some words about the tests}\smallskip
As always, the hard part is coming up with meaningful abbreviations for the operations
\par

'r' for 'reference counter' is a good enough mnemonic. that'll be the prefix.
\par

verbs are: add, delete, find, increase, decrease, hold, unhold, collect, get. getlen and getcount are the helpers for tests.
\par

ri for init. ra for 'add'. 'rm' for delete. 'ru' for update counter, which can be negative and positive, 'rf' for find. 'rh' for hold/unhold (with arg), 'rg' for get. 'rx' will be an aux function, which will implement getlen() and getcount() based on an input argument.
\par

punching these into the glossary with a TODO.
\par

Going to start writing words for the tests now.
\par

I got words and initial unit tests set up for add/rm, find, and incr/decr. I shall return for the rest later.
\par

{\smallskip\noindent\bf\font\titlefont=cmr12 at 13pt\titlefont $\bullet$ 20:14: finish writing up initial unit tests}\smallskip
Huh, forgot about collect. 'rc' is taken, so 'rs' for 'sweep' should work.
\par

{\smallskip\noindent\bf\font\titlefont=cmr12 at 13pt\titlefont $\bullet$ 21:05: coding some things up}\smallskip
Just some initial scaffolding for C functions and corresponding words.
\par

{\smallskip\noindent\bf\font\titlefont=cmr12 at 13pt\titlefont $\bullet$ 21:52: finish up implementation}\smallskip
WHAT. how is it almost 23:00. This is why it's good to cut back on coding. You lose so much time in your life.
\par

This of course is the hour to get a weird by with the get command. Why is it that the final number of active slots is 4 and not 3?
\par

okay, fixed it. I forgot to check for only active entries (aka entries with non-zero addresses).
\par

I think these new abbreviations triggered a bug for some old test code. {\tt test\_rw} is using a command it calls {\tt clr}, when I think it's supposed to be {\tt cl}.
\par

I think we are done here.
\par

\xrdef{node35}
{\medskip\noindent\bf\font\titlefont=cmss17 at 20pt\titlefont ugen stack}\medskip
{\medskip\noindent\bf\font\titlefont=cmitt12 at 17pt\titlefont 2026-01-31}\medskip
{\smallskip\noindent\bf\font\titlefont=cmr12 at 13pt\titlefont $\bullet$ 14:34: creating new tasks for ugens, writing words}\smallskip
{\smallskip\noindent\bf\font\titlefont=cmr12 at 13pt\titlefont $\bullet$ 14:50: initial words}\smallskip
A stack is used to share data between ugens. When a unit generator is instantiated, it reads inputs by popping data from the stack, and writes output values by pushing data onto the stack. This, combined with a reference counting system, builds up an abstraction that is able to recycle blocks within the scope of a patch (a collection of unit generators).
\par

{\smallskip\noindent\bf\font\titlefont=cmr12 at 13pt\titlefont $\bullet$ 14:55: Garbage collection after ports, order matters}\smallskip
I may have mentioned this already. But it's one of those things I can see slipping from my mind. So, we're writing it down again.
\par

Unit generators need to be able to share and reuse blocks. A stack and a reference counting system is used to help with this.
\par

A ugen has a certain number of input ports and output ports. A ugen generator will allocate data required for the DSP state, set up the ports, and then finally do garbage collection on the reference counter.
\par

A pop operation decreases reference count associated with a particular memory address of a block. A push operation will increase it. A garbage collect operation will mark that block as available for re-use.
\par

The garbage collection happens last, after all stack operations. Any time a new output parameter is pushed onto the stack, it will be unable to claim a recently popped block used by the input, guaranteeing that input/output blocks will be disjoing. Making the output/input blocks disjoint prevents things like accidental overwrites to happen.
\par

{\smallskip\noindent\bf\font\titlefont=cmr12 at 13pt\titlefont $\bullet$ 15:19: packing the stack and refcounter together using some arithmetic}\smallskip
One refcount entry needs 18 bits of data. 64 slots is 1152 bits of data, conveniently 36 words of data. 28 words left for the stack.
\par

slot indexes need 6 bits of data. If the stack uses an array, it's possible to have upper limit of approximiately 149 for a stack size. Round to the nearest power of 2, 128. 768 bits, 24 words. But, that doesn't include the stack pointer, which would need to be at least 8 bits to include 0.
\par

Alternatively, use an array structure, and have the first entry contain the stack pointer (the number of items on the stack). a 6-bit number would have a range 0-63, or a stack size of 63. a 7-bit number would have a range 0-127, or a stack size of 127.
\par

a 6-bit stack needs an array of 64 6-bit numbers, 12 words. A 7-bit stack needs an array of 128 7-bit numbers, 28 words.
\par

It turns out that 64 18-bit numbers plus 128 7-bit numbers yields 32 64-bit numbers, or one block. How lovely.
\par

{\medskip\noindent\bf\font\titlefont=cmitt12 at 17pt\titlefont 2026-02-14}\medskip
{\smallskip\noindent\bf\font\titlefont=cmr12 at 13pt\titlefont $\bullet$ 10:01: hoping to do some initial work today on this}\smallskip
This stack is designed to work with the reference counter structure, so it's all on the same block.
\par

Push/pop operations should implicitly update refcounts.
\par

I was currently envisioning 7-bit data types. This is good enough for addresses relative to the block. But, it does not handle other data types like constants. This would be desirable right?
\par

Looks like there's 640 bits to play with after subtracting the size of the refcount structure. That's 20 words, so a stack size of 19. Or it could be 40 half-words, with a stack size of 39.
\par

I think I'm leaning towards a smaller stack, larger data type. I get the feeling I'm going to be adding more types of data before I run out of stack space. 32 bits, with the first two bits being a type, 30 bits for data. Stack size of 19. 2 types (constant or block) will be used right now, with the 2 other types TBD. Constants can be encoded with 15.15 signed data encoding or something and that should be good enough.
\par

{\smallskip\noindent\bf\font\titlefont=cmr12 at 13pt\titlefont $\bullet$ 13:40: write up initial unit tests}\smallskip
First, let's take a look at the commands I'll need to implement, and their names.
\par

Operations include: push/pop, hold/unhold (implicit), dup, and rot. There are other stack operations in forths, so having a good single letter namespace would be nice.
\par

I've shifted from calling this a "ugen stack" to a "parameter stack", because it more precisely describes the data being handled, which is any valid parameter for a unit generator.
\par

The 'p' prefix is reasonably open. pr, pa, pw, and pe are taken.
\par

po: pop. pu: push. pi: init stack (give it block). ph: hold/unhold (this can be given a parameter). pd: dup. pR: rotate (capital letters work, as long as it's not A-F).
\par

should stacks implicitly add addresses to the RC? I'm going to say... no. That's just one extra thing that can go wrong. In practice, I imagine a pool of blocks will be initialized and added to the RC pool. Then, there'd be operations that select available blocks, abstracting away the need to know the blocks of memory addresses.
\par

Pool is TBD. adding a task for it, with some words.
\par

'px' I'll use for pstack aux functions. that was a nice thing to have with the RC stuff.
\par

Just remembered that stacks are going to need types. Gotta update the push commands to take in a type flag so it knows what to do with the bits.
\par

\xrdef{node36}
{\medskip\noindent\bf\font\titlefont=cmss17 at 20pt\titlefont bttyp}\medskip
{\medskip\noindent\bf\font\titlefont=cmitt12 at 17pt\titlefont 2026-01-31}\medskip
{\smallskip\noindent\bf\font\titlefont=cmr12 at 13pt\titlefont $\bullet$ 20:39: bttyp initial coding}\smallskip
I've had this idea for a 1-bit typesetting system using mac fonts for years now. Going to see what I can accomplish tonight.
\par

I have some macfont parsing and rendering code somewhere which does most of the hard bits. I'll want to adapt it a bit to fit with the framework I have in mind.
\par

There's a serial language which I'll make a parser for. And I'll also make a human-readable version that assembles to the parser (in AWK).
\par

Setting up some initial scaffolding. Something that writes a whole page to disk.
\par

Did some fiddling, and decided a page should be 512x832. It's a rectangle close to the golden ratio, while also being a multiple of 8 both ways.
\par

copying over files from btprnt, and chicago12
\par

I have a blank PNG file now.
\par

I have an 'A' rendered now.
\par

Working on an initial bytestream interpreter. I figured out how to get Awk to print special characters. Now I'll feed that into my program.
\par

Some initial text primitives made. I have a few lines of text being rendered in chicago12. Not bad for a first start.
\par

{\medskip\noindent\bf\font\titlefont=cmitt12 at 17pt\titlefont 2026-02-03}\medskip
{\smallskip\noindent\bf\font\titlefont=cmr12 at 13pt\titlefont $\bullet$ 20:25: bttyp: abstraction barriers}\smallskip
thinking about the layers of abstraction required to make bttyp to be useful. The current prototype is a proof of concept that glues together existing code I once wrote, but there needs to be more in order to make documents.
\par

In this brainstorm, I draw a line between in the software abstractions. low-level "from below" stuff deals with rasterization, while high-level "from above" stuff deals with concepts like alignment, word breaking, etc. The "from above" stuff needs to be built and thought on more, and will also influence where the "from below" stuff will be headed.
\par

\xrdef{node37}
{\medskip\noindent\bf\font\titlefont=cmss17 at 20pt\titlefont block pool}\medskip
{\medskip\noindent\bf\font\titlefont=cmitt12 at 17pt\titlefont 2026-02-14}\medskip
{\smallskip\noindent\bf\font\titlefont=cmr12 at 13pt\titlefont $\bullet$ 13:57: A block pool abstraction would be nice to have}\smallskip
This would abstract away the tediousness of allocating blocks from the global block stack and keeping track of memory addresses manually.
\par

A new block is needed any time an output block is created for a ugen. So steps would be: set up input blocks, set up output block(s), sweep.
\par

If I were making, say a sine ugen with frequency and amplitude as args, the steps would be: set up ugen with 2 ins, 1 out. pop amp, assign to port 0, pop freq, assign to port 1, new block for output, assign to port 2. alloc data, initialize, set to ugen. sweep.
\par

{\smallskip\noindent\bf\font\titlefont=cmr12 at 13pt\titlefont $\bullet$ 19:23: implementation}\smallskip
\bye
